\documentclass[11pt,a4paper]{article}
\usepackage[utf8]{inputenc}
\usepackage[T1]{fontenc}
\usepackage[french]{babel}
\usepackage{amsmath,amssymb,amsfonts,amsthm,mathrsfs}
\usepackage[margin=2cm]{geometry}
\usepackage{enumitem}
\usepackage{xcolor}
\usepackage{tcolorbox}
\usepackage{hyperref}

\tcbuselibrary{skins,breakable}

\tcolorboxenvironment{proof}{
  blanker,breakable,left=3mm,
  before skip=10pt,after skip=10pt,
  borderline west={1.5mm}{0pt}{gray!40}
}

\newtcolorbox{correctionbox}{
  colback=blue!5!white,
  colframe=blue!60!black,
  title=\textbf{Correction détaillée},
  fonttitle=\bfseries,
  arc=2mm,
  breakable
}

\newcommand{\N}{\mathbb{N}}
\newcommand{\Z}{\mathbb{Z}}
\newcommand{\Q}{\mathbb{Q}}
\newcommand{\R}{\mathbb{R}}
\newcommand{\C}{\mathbb{C}}
\providecommand{\euro}{\text{\officialeuro}}

\title{\textbf{Chapitre 2 : Le Second Degré et Polynômes}\\\large Fiche complète de TD d'approfondissement (Exercices 1 à 41) et corrections détaillées\\{\small Niveau : Première Spécialité Mathématiques (Lycée Louis-le-Grand)}}
\author{Soutien Scolaire -- Mathématiques}
\date{Année scolaire 2021--2022}

\begin{document}

\maketitle

\tableofcontents
\vspace{1cm}
\hrule
\vspace{1cm}

\section*{Première partie : Exercices 1 à 20}

% ----------------------------------------------------
% EXERCICE 1
% ----------------------------------------------------
\subsection*{Exercice 1 : Un résultat de cours}
Démontrer que pour tous nombres réels $x$ et $y$ :
\[
x^2 = y^2 \iff (x = y \quad \text{ou} \quad x = -y).
\]

\begin{correctionbox}
Soient $x$ et $y$ deux nombres réels.
\begin{enumerate}
    \item Par soustraction membre à membre :
    \[
    x^2 = y^2 \iff x^2 - y^2 = 0.
    \]
    \item On factorise le membre de gauche à l'aide de l'identité remarquable $a^2 - b^2 = (a-b)(a+b)$ :
    \[
    x^2 - y^2 = 0 \iff (x - y)(x + y) = 0.
    \]
    \item Dans le corps des réels $\R$ (qui est intègre), un produit de facteurs est nul si et seulement si l'un au moins des facteurs est nul (règle du produit nul) :
    \[
    (x - y)(x + y) = 0 \iff x - y = 0 \quad \text{ou} \quad x + y = 0 \iff x = y \quad \text{ou} \quad x = -y.
    \]
\end{enumerate}
Chaque étape étant une équivalence logique stricte, on en déduit que :
\[
\forall (x, y) \in \R^2, \quad x^2 = y^2 \iff (x = y \quad \text{ou} \quad x = -y).
\]

\textbf{Remarque :} En utilisant la fonction valeur absolue, ce résultat fondamental s'énonce de façon équivalente sous la forme :
\[
x^2 = y^2 \iff |x| = |y|.
\]
Il met en garde contre l'erreur fréquente consistant à « prendre la racine carrée » sans précaution, puisque $\sqrt{x^2} = |x|$ et non $x$.
\end{correctionbox}

% ----------------------------------------------------
% EXERCICE 2
% ----------------------------------------------------
\subsection*{Exercice 2 : Équation rationnelle et valeur interdite}
Résoudre dans $\R$ l'équation :
\[
\frac{x^2 + 2x + 1}{x + 1} = 0.
\]

\begin{correctionbox}
\begin{enumerate}
    \item \textbf{Ensemble de définition (valeur interdite) :}
    Le quotient $\frac{x^2+2x+1}{x+1}$ est défini si et seulement si son dénominateur est non nul :
    \[
    x + 1 \neq 0 \iff x \neq -1.
    \]
    L'ensemble de validité de l'équation est donc $\mathcal{D} = \R \setminus \{-1\}$.

    \item \textbf{Résolution sur le domaine $\mathcal{D}$ :}
    Pour tout $x \in \mathcal{D}$, un quotient est nul si et seulement si son numérateur s'annule :
    \[
    \frac{x^2 + 2x + 1}{x + 1} = 0 \iff x^2 + 2x + 1 = 0.
    \]
    On reconnaît l'identité remarquable du carré d'une somme :
    \[
    (x + 1)^2 = 0 \iff x + 1 = 0 \iff x = -1.
    \]

    \item \textbf{Confrontation avec le domaine de définition :}
    L'unique solution candidate est $x = -1$.
    Or, $-1 \notin \mathcal{D}$ car il annule le dénominateur (il s'agit d'une valeur interdite).

    \item \textbf{Conclusion :}
    L'équation ne possède aucune solution dans $\R$ :
    \[
    \mathcal{S} = \emptyset.
    \]
\end{enumerate}
\textbf{Remarque méthodologique :} Simplifier directement l'expression par $x+1$ sans poser le domaine préalable :
\[
\frac{(x+1)^2}{x+1} = x+1
\]
induirait en erreur et mènerait à la fausse conclusion $x = -1$. Une simplification de fraction n'est valide que sur son ensemble de définition !
\end{correctionbox}

% ----------------------------------------------------
% EXERCICE 3
% ----------------------------------------------------
\subsection*{Exercice 3 : Somme d'inverses et réduction au même dénominateur}
Résoudre dans $\R$ l'équation :
\[
\frac{1}{x-1} + \frac{1}{x-2} + \frac{1}{x-3} = 0.
\]

\begin{correctionbox}
\begin{enumerate}
    \item \textbf{Domaine de définition :}
    L'équation a un sens si et seulement si aucun des dénominateurs ne s'annule :
    \[
    x - 1 \neq 0, \quad x - 2 \neq 0 \quad \text{et} \quad x - 3 \neq 0.
    \]
    Le domaine de résolution est donc :
    \[
    \mathcal{D} = \R \setminus \{1\,,\, 2\,,\, 3\}.
    \]

    \item \textbf{Réduction au même dénominateur :}
    Pour tout $x \in \mathcal{D}$, le dénominateur commun est le produit $(x-1)(x-2)(x-3)$.
    L'équation devient :
    \[
    \frac{(x-2)(x-3) + (x-1)(x-3) + (x-1)(x-2)}{(x-1)(x-2)(x-3)} = 0.
    \]
    Pour $x \in \mathcal{D}$, cette équation équivaut à l'annulation du numérateur $N(x) = 0$.
    Développons séparément chacun des produits :
    \begin{align*}
    (x-2)(x-3) &= x^2 - 5x + 6, \\
    (x-1)(x-3) &= x^2 - 4x + 3, \\
    (x-1)(x-2) &= x^2 - 3x + 2.
    \end{align*}
    En additionnant ces trois polynômes :
    \[
    N(x) = (x^2 - 5x + 6) + (x^2 - 4x + 3) + (x^2 - 3x + 2) = 3x^2 - 12x + 11.
    \]
    L'équation se ramène donc, pour $x \in \mathcal{D}$, à :
    \[
    3x^2 - 12x + 11 = 0.
    \]

    \item \textbf{Résolution de l'équation du second degré :}
    Calculons le discriminant du trinôme :
    \[
    \Delta = (-12)^2 - 4 \times 3 \times 11 = 144 - 132 = 12 > 0.
    \]
    Comme $\Delta > 0$, l'équation admet deux racines réelles distinctes :
    \[
    x_1 = \frac{-(-12) - \sqrt{12}}{2 \times 3} = \frac{12 - 2\sqrt{3}}{6} = \frac{6 - \sqrt{3}}{3} = 2 - \frac{\sqrt{3}}{3},
    \]
    \[
    x_2 = \frac{-(-12) + \sqrt{12}}{2 \times 3} = \frac{12 + 2\sqrt{3}}{6} = \frac{6 + \sqrt{3}}{3} = 2 + \frac{\sqrt{3}}{3}.
    \]

    \item \textbf{Vérification d'appartenance au domaine $\mathcal{D}$ :}
    Puisque $1 < \sqrt{3} < 2$, on a $\frac{1}{3} < \frac{\sqrt{3}}{3} < \frac{2}{3} < 1$.
    Par suite :
    \[
    1 < \frac{4}{3} < 2 - \frac{\sqrt{3}}{3} < \frac{5}{3} < 2 \implies x_1 \notin \{1, 2, 3\},
    \]
    \[
    2 < \frac{7}{3} < 2 + \frac{\sqrt{3}}{3} < \frac{8}{3} < 3 \implies x_2 \notin \{1, 2, 3\}.
    \]
    Les deux solutions candidates appartiennent bien à l'ensemble de validité $\mathcal{D}$.

    \item \textbf{Conclusion :}
    L'ensemble des solutions de l'équation est :
    \[
    \mathcal{S} = \left\{ \frac{6 - \sqrt{3}}{3}\,,\, \frac{6 + \sqrt{3}}{3} \right\} = \left\{ 2 - \frac{\sqrt{3}}{3}\,,\, 2 + \frac{\sqrt{3}}{3} \right\}.
    \]
\end{enumerate}
\end{correctionbox}

% ----------------------------------------------------
% EXERCICE 4
% ----------------------------------------------------
\subsection*{Exercice 4 : Résolution sans discriminant (Relations de Viète)}
Résoudre dans $\R$, sans calculer le discriminant $\Delta$, l'équation :
\[
2021x^2 + x - 2022 = 0.
\]

\begin{correctionbox}
Soit $P(x) = 2021x^2 + x - 2022$. Il s'agit d'un trinôme du second degré $ax^2 + bx + c$ avec $a = 2021$, $b = 1$ et $c = -2022$.
\begin{enumerate}
    \item \textbf{Recherche d'une racine évidente par la somme des coefficients :}
    Calculons la somme des coefficients de ce trinôme :
    \[
    a + b + c = 2021 + 1 + (-2022) = 2022 - 2022 = 0.
    \]
    Or, pour tout polynôme $P(x) = ax^2 + bx + c$, on a $P(1) = a(1)^2 + b(1) + c = a + b + c$.
    Comme la somme des coefficients est nulle, le nombre $1$ est racine évidente de l'équation :
    \[
    P(1) = 0 \implies x_1 = 1.
    \]

    \item \textbf{Détermination de la seconde racine via le produit des racines :}
    Soit $x_2$ la seconde racine de l'équation. D'après les relations de Viète liant les racines aux coefficients d'un polynôme du second degré :
    \[
    P = x_1 \times x_2 = \frac{c}{a}.
    \]
    En remplaçant par les valeurs numériques et en utilisant $x_1 = 1$ :
    \[
    1 \times x_2 = \frac{-2022}{2021} \implies x_2 = -\frac{2022}{2021}.
    \]

    \item \textbf{Vérification par la somme des racines :}
    On contrôle la cohérence à l'aide de la somme des racines :
    \[
    S = x_1 + x_2 = 1 + \left(-\frac{2022}{2021}\right) = \frac{2021 - 2022}{2021} = -\frac{1}{2021} = -\frac{b}{a},
    \]
    ce qui corrobore parfaitement le résultat.

    \item \textbf{Conclusion :}
    L'ensemble des solutions de l'équation est :
    \[
    \mathcal{S} = \left\{ -\frac{2022}{2021}\,,\, 1 \right\}.
    \]
\end{enumerate}
\end{correctionbox}

% ----------------------------------------------------
% EXERCICE 5
% ----------------------------------------------------
\subsection*{Exercice 5 : Problème géométrique d'optimisation d'aire}
Soit $ABCD$ un carré de côté $4\text{ cm}$. On choisit un point $E$ sur le segment $[AB]$.
On construit à l'intérieur du carré :
\begin{itemize}
    \item le carré $AEFG$ de côté $AE$ ;
    \item le triangle $EHB$, rectangle et isocèle en $H$, d'hypoténuse $[EB]$ portée par le segment $[AB]$.
\end{itemize}
Déterminer la position du point $E$ sur le segment $[AB]$ pour que la somme des aires du carré $AEFG$ et du triangle $EHB$ soit égale à $5\text{ cm}^2$.

\begin{correctionbox}
\begin{enumerate}
    \item \textbf{Choix de l'inconnue et domaine géométrique :}
    On pose $x = AE$ en centimètres.
    Le point $E$ appartenant au segment $[AB]$ de longueur $4\text{ cm}$, l'inconnue $x$ vérifie :
    \[
    x \in [0\,,\, 4].
    \]

    \item \textbf{Calcul des aires en fonction de $x$ :}
    \begin{itemize}
        \item \textbf{Aire du carré $AEFG$ :}
        Le côté du carré est $AE = x$, donc son aire est :
        \[
        \mathcal{A}_1(x) = x^2.
        \]
        \item \textbf{Aire du triangle $EHB$ :}
        Puisque $E \in [AB]$, la longueur de l'hypoténuse $[EB]$ est $EB = AB - AE = 4 - x$.
        Le triangle $EHB$ étant rectangle et isocèle en $H$, on a $EH = HB$. D'après le théorème de Pythagore dans le triangle $EHB$ :
        \[
        EH^2 + HB^2 = EB^2 \iff 2 EH^2 = (4-x)^2 \iff EH^2 = \frac{(4-x)^2}{2}.
        \]
        L'aire du triangle rectangle $EHB$ est donnée par la moitié du produit des côtés de l'angle droit :
        \[
        \mathcal{A}_2(x) = \frac{1}{2} \times EH \times HB = \frac{1}{2} EH^2 = \frac{1}{2} \times \frac{(4-x)^2}{2} = \frac{(4-x)^2}{4}.
        \]
    \end{itemize}
    L'aire totale colorée de la figure est donc :
    \[
    \mathcal{A}(x) = \mathcal{A}_1(x) + \mathcal{A}_2(x) = x^2 + \frac{(4-x)^2}{4}.
    \]

    \item \textbf{Mise en équation :}
    La condition imposée est $\mathcal{A}(x) = 5$, ce qui se traduit par :
    \[
    x^2 + \frac{(4-x)^2}{4} = 5.
    \]
    En multipliant les deux membres par $4$ :
    \begin{align*}
    4x^2 + (4-x)^2 = 20 &\iff 4x^2 + (16 - 8x + x^2) = 20 \\
    &\iff 5x^2 - 8x + 16 - 20 = 0 \\
    &\iff 5x^2 - 8x - 4 = 0.
    \end{align*}

    \item \textbf{Résolution de l'équation du second degré :}
    Calculons le discriminant :
    \[
    \Delta = (-8)^2 - 4 \times 5 \times (-4) = 64 + 80 = 144 = 12^2 > 0.
    \]
    Les deux racines réelles sont :
    \[
    x_1 = \frac{-(-8) - 12}{2 \times 5} = \frac{8 - 12}{10} = -\frac{4}{10} = -0{,}4,
    \]
    \[
    x_2 = \frac{-(-8) + 12}{2 \times 5} = \frac{8 + 12}{10} = \frac{20}{10} = 2.
    \]

    \item \textbf{Confrontation avec la contrainte géométrique :}
    Puisque $x$ représente une distance sur le segment $[AB]$, on doit avoir $x \in [0, 4]$.
    La valeur $x_1 = -0{,}4$ est strictement négative, elle est donc exclue.
    La seule valeur admissible est $x_2 = 2$.

    \item \textbf{Conclusion :}
    L'aire demandée est égale à $5\text{ cm}^2$ si et seulement si $AE = 2\text{ cm}$, ce qui signifie que le point $E$ est exactement le \textbf{milieu} du segment $[AB]$.
\end{enumerate}
\end{correctionbox}

% ----------------------------------------------------
% EXERCICE 6
% ----------------------------------------------------
\subsection*{Exercice 6 : Inégalité fondamentale sur la somme d'un réel et de son inverse}
\begin{enumerate}
    \item Démontrer que pour tout nombre réel $x > 0$ :
    \[
    x + \frac{1}{x} \ge 2.
    \]
    \item Déterminer la condition nécessaire et suffisante pour que l'égalité ait lieu.
    \item Établir l'inégalité analogue satisfaite pour tout réel $x < 0$.
\end{enumerate}

\begin{correctionbox}
\begin{enumerate}
    \item Soit $x \in \R_+^* = ]0\,,\, +\infty[$.
    Pour comparer la quantité $x + \frac{1}{x}$ au nombre $2$, étudions le signe de leur différence :
    \[
    D(x) = x + \frac{1}{x} - 2.
    \]
    Réduisons au même dénominateur $x > 0$ :
    \[
    D(x) = \frac{x^2 - 2x + 1}{x} = \frac{(x - 1)^2}{x}.
    \]
    Or :
    \begin{itemize}
        \item Pour tout réel $x$, $(x - 1)^2 \ge 0$ car le carré d'un nombre réel est toujours positif ou nul ;
        \item Par hypothèse, $x > 0$.
    \end{itemize}
    Le quotient d'un nombre positif ou nul par un nombre strictement positif est positif ou nul :
    \[
    \forall x > 0, \quad \frac{(x-1)^2}{x} \ge 0 \iff x + \frac{1}{x} - 2 \ge 0 \iff x + \frac{1}{x} \ge 2.
    \]

    \item L'égalité a lieu si et seulement si la différence s'annule :
    \[
    x + \frac{1}{x} = 2 \iff \frac{(x-1)^2}{x} = 0 \iff (x-1)^2 = 0 \iff x - 1 = 0 \iff x = 1.
    \]
    Le minimum de la fonction $x \mapsto x + \frac{1}{x}$ sur $]0, +\infty[$ est donc atteint en un unique point : $x = 1$, et vaut $2$.

    \item Supposons maintenant $x < 0$. Alors $-x > 0$.
    Appliquons le résultat de la première question au réel strictement positif $y = -x > 0$ :
    \[
    (-x) + \frac{1}{-x} \ge 2 \iff -\left(x + \frac{1}{x}\right) \ge 2.
    \]
    En multipliant les deux membres par $-1$ (ce qui inverse le sens de l'inégalité) :
    \[
    \forall x < 0, \quad x + \frac{1}{x} \le -2.
    \]
    L'égalité a lieu si et seulement si $-x = 1 \iff x = -1$.
\end{enumerate}
\textbf{Remarque culturelle :} Ce résultat est une illustration immédiate de l'inégalité arithmético-géométrique (AM-GM) appliquée aux deux réels strictement positifs $a = x$ et $b = \frac{1}{x}$ :
\[
\frac{a + b}{2} \ge \sqrt{ab} \iff \frac{x + \frac{1}{x}}{2} \ge \sqrt{x \times \frac{1}{x}} = 1 \iff x + \frac{1}{x} \ge 2.
\]
\end{correctionbox}

% ----------------------------------------------------
% EXERCICE 7
% ----------------------------------------------------
\subsection*{Exercice 7 : Modélisation d'intérêts composés (Mathématiques financières)}
Un épargnant dépose chaque année la somme de $1\,000$~\euro{} sur un compte bancaire rémunéré à un taux annuel constant à intérêts composés noté $t = \frac{p}{100}$ :
\begin{itemize}
    \item un premier versement de $1\,000$~\euro{} à la date initiale ($T = 0$) ;
    \item un deuxième versement de $1\,000$~\euro{} un an plus tard ($T = 1$) ;
    \item un troisième versement de $1\,000$~\euro{} deux ans après l'ouverture ($T = 2$).
\end{itemize}
Immédiatement après le troisième versement, le capital total accumulé s'élève à $3\,090{,}90$~\euro{}.
Déterminer la valeur exacte du taux d'intérêt $t$, puis le pourcentage annuel $p$.

\begin{correctionbox}
\begin{enumerate}
    \item \textbf{Modélisation de la valeur acquise au bout de 2 ans :}
    À la date $T = 2$ :
    \begin{itemize}
        \item Le premier versement de $1\,000$~\euro{} a été placé pendant $2$ ans ; avec le taux d'intérêt annuel composé $t$, il a été multiplié par le coefficient $(1+t)^2$ et vaut $1\,000(1+t)^2$.
        \item Le deuxième versement de $1\,000$~\euro{} a été placé pendant $1$ an ; il vaut $1\,000(1+t)$.
        \item Le troisième versement de $1\,000$~\euro{} vient d'être déposé à $T = 2$ ; sa valeur est donc $1\,000$.
    \end{itemize}
    La somme totale disponible s'écrit donc :
    \[
    C(t) = 1\,000(1+t)^2 + 1\,000(1+t) + 1\,000 = 3\,090{,}90.
    \]

    \item \textbf{Mise sous forme d'équation polynomiale du second degré :}
    En factorisant par $1\,000$ :
    \[
    1\,000 \bigl( (1+t)^2 + (1+t) + 1 \bigr) = 3\,090{,}90.
    \]
    Développons le terme entre parenthèses :
    \[
    (1+t)^2 + (1+t) + 1 = (1 + 2t + t^2) + 1 + t + 1 = t^2 + 3t + 3.
    \]
    L'équation devient :
    \[
    1\,000(t^2 + 3t + 3) = 3\,090{,}90 \iff t^2 + 3t + 3 = 3{,}0909 \iff t^2 + 3t - 0{,}0909 = 0.
    \]

    \item \textbf{Résolution de l'équation :}
    Calculons le discriminant :
    \[
    \Delta = 3^2 - 4 \times 1 \times (-0{,}0909) = 9 + 0{,}3636 = 9{,}3636.
    \]
    Cherchons la racine carrée de ce discriminant :
    \[
    9{,}3636 = \frac{93\,636}{10\,000} = \left(\frac{306}{100}\right)^2 = (3{,}06)^2.
    \]
    En effet : $306^2 = (300 + 6)^2 = 90\,000 + 3\,600 + 36 = 93\,636$.
    Ainsi, $\sqrt{\Delta} = 3{,}06$. Les deux racines sont :
    \[
    t_1 = \frac{-3 - 3{,}06}{2} = \frac{-6{,}06}{2} = -3{,}03,
    \]
    \[
    t_2 = \frac{-3 + 3{,}06}{2} = \frac{0{,}06}{2} = 0{,}03.
    \]

    \item \textbf{Conclusion :}
    Un taux de placement bancaire étant nécessairement positif ($t > 0$), la racine négative $t_1 = -3{,}03$ est dépourvue de sens économique.
    On en déduit que le taux d'intérêt est :
    \[
    t = 0{,}03 \quad \text{soit un taux annuel de } p = 3\,\%.
    \]
\end{enumerate}
\end{correctionbox}

% ----------------------------------------------------
% EXERCICE 8
% ----------------------------------------------------
\subsection*{Exercice 8 : Détermination de l'expression d'un trinôme par interpolation}
Déterminer l'expression développée et réduite de la fonction polynôme du second degré $f$ dont la courbe représentative (une parabole $\mathcal{P}$) passe par les points de coordonnées $(-6\,;\,0)$, $(1\,;\,0)$ et $(0\,;\,2)$.

\begin{correctionbox}
Soit $f(x) = ax^2 + bx + c$ avec $a \in \R^*$ le trinôme du second degré recherché.
\begin{enumerate}
    \item \textbf{Utilisation des racines et de la forme factorisée :}
    La parabole coupe l'axe des abscisses aux points d'abscisses $-6$ et $1$. Les nombres $-6$ et $1$ sont donc les racines de $f$ :
    \[
    f(-6) = 0 \quad \text{et} \quad f(1) = 0.
    \]
    D'après le cours, tout trinôme admettant deux racines réelles distinctes $x_1$ et $x_2$ s'écrit sous forme factorisée :
    \[
    f(x) = a(x - x_1)(x - x_2) = a\bigl(x - (-6)\bigr)(x - 1) = a(x + 6)(x - 1).
    \]

    \item \textbf{Détermination du coefficient dominant $a$ :}
    La parabole passe par le point $(0\,;\,2)$, ce qui signifie que l'ordonnée à l'origine vaut $2$, soit $f(0) = 2$ :
    \[
    f(0) = 2 \iff a(0 + 6)(0 - 1) = 2 \iff -6a = 2 \iff a = -\frac{2}{6} = -\frac{1}{3}.
    \]

    \item \textbf{Développement et forme réduite :}
    On développe l'expression obtenue :
    \[
    f(x) = -\frac{1}{3}(x + 6)(x - 1) = -\frac{1}{3}(x^2 + 5x - 6) = -\frac{1}{3}x^2 - \frac{5}{3}x + 2.
    \]

    \item \textbf{Vérification des éléments géométriques :}
    \begin{itemize}
        \item $f(0) = 2$ ;
        \item $f(1) = -\frac{1}{3} - \frac{5}{3} + 2 = -2 + 2 = 0$ ;
        \item $f(-6) = -\frac{1}{3}(36) - \frac{5}{3}(-6) + 2 = -12 + 10 + 2 = 0$.
    \end{itemize}
    Toutes les conditions de l'énoncé sont rigoureusement vérifiées.
\end{enumerate}
\textbf{Conclusion :} L'expression recherchée est :
\[
f(x) = -\frac{1}{3}x^2 - \frac{5}{3}x + 2.
\]
\end{correctionbox}

% ----------------------------------------------------
% EXERCICE 9
% ----------------------------------------------------
\subsection*{Exercice 9 : Équation avec racine carrée et changement de variable}
Résoudre dans $\R$ l'équation :
\[
6x - 13\sqrt{x} - 5 = 0.
\]

\begin{correctionbox}
\begin{enumerate}
    \item \textbf{Ensemble d'admissibilité :}
    L'expression $\sqrt{x}$ n'a de sens réel que si le radicande est positif ou nul.
    Le domaine de validité est donc :
    \[
    \mathcal{D} = [0\,,\, +\infty[.
    \]

    \item \textbf{Changement de variable :}
    Pour tout $x \in \mathcal{D}$, on a $x = \bigl(\sqrt{x}\bigr)^2$.
    Posons la variable auxiliaire $X = \sqrt{x}$.
    Une racine carrée étant toujours positive ou nulle, la nouvelle inconnue $X$ doit vérifier la contrainte :
    \[
    X \ge 0.
    \]
    Sous cette condition, l'équation se réécrit comme une équation polynomiale du second degré en $X$ :
    \[
    6X^2 - 13X - 5 = 0.
    \]

    \item \textbf{Résolution en $X$ :}
    Calculons le discriminant :
    \[
    \Delta = (-13)^2 - 4 \times 6 \times (-5) = 169 + 120 = 289.
    \]
    Puisque $289 = 17^2 > 0$, l'équation en $X$ admet deux racines réelles distinctes :
    \[
    X_1 = \frac{-(-13) - 17}{2 \times 6} = \frac{13 - 17}{12} = -\frac{4}{12} = -\frac{1}{3},
    \]
    \[
    X_2 = \frac{-(-13) + 17}{2 \times 6} = \frac{13 + 17}{12} = \frac{30}{12} = \frac{5}{2}.
    \]

    \item \textbf{Retour à la variable $x$ :}
    \begin{itemize}
        \item La valeur candidate $X_1 = -\frac{1}{3}$ est strictement négative, ce qui contredit la condition impérative $X = \sqrt{x} \ge 0$. Cette racine est donc \textbf{rejetée}.
        \item La valeur candidate $X_2 = \frac{5}{2}$ est strictement positive, donc admissible :
        \[
        \sqrt{x} = \frac{5}{2} \iff x = \left(\frac{5}{2}\right)^2 = \frac{25}{4}.
        \]
    \end{itemize}

    \item \textbf{Conclusion :}
    L'équation admet une unique solution réelle :
    \[
    \mathcal{S} = \left\{ \frac{25}{4} \right\}.
    \]
\end{enumerate}
\end{correctionbox}

% ----------------------------------------------------
% EXERCICE 10
% ----------------------------------------------------
\subsection*{Exercice 10 : Inéquations bicarrées}
Résoudre dans $\R$ les inéquations suivantes :
\begin{enumerate}
    \item $x^4 + 13x^2 + 36 > 0$
    \item $x^4 - 13x^2 + 36 > 0$
\end{enumerate}

\begin{correctionbox}
\begin{enumerate}
    \item \textbf{Résolution de l'inéquation (1) :} $x^4 + 13x^2 + 36 > 0$.
    Pour tout nombre réel $x \in \R$ :
    \[
    x^2 \ge 0 \implies x^4 = (x^2)^2 \ge 0 \quad \text{et} \quad 13x^2 \ge 0.
    \]
    Comme $36 > 0$, la somme d'éléments positifs ou nuls et d'une constante strictement positive est strictement positive :
    \[
    \forall x \in \R, \quad x^4 + 13x^2 + 36 \ge 36 > 0.
    \]
    L'inégalité est vérifiée pour tout nombre réel. Par conséquent :
    \[
    \mathcal{S}_1 = \R.
    \]

    \item \textbf{Résolution de l'inéquation (2) :} $x^4 - 13x^2 + 36 > 0$.
    \begin{itemize}
        \item \textbf{Changement de variable :} Posons $X = x^2$. L'inconnue $X$ est soumise à la condition $X \ge 0$.
        Le polynôme s'écrit sous la forme du trinôme :
        \[
        P(X) = X^2 - 13X + 36.
        \]
        \item \textbf{Racines et factorisation de $P(X)$ :}
        On cherche deux nombres ayant pour somme $S = 13$ et pour produit $P = 36$. Ces deux nombres sont immédiatement $4$ et $9$.
        Le trinôme se factorise donc sous la forme :
        \[
        P(X) = (X - 4)(X - 9).
        \]
        \item \textbf{Signe du trinôme en $X$ :}
        Le coefficient dominant de $P(X)$ est $1 > 0$. Le trinôme est strictement positif à l'extérieur de ses racines :
        \[
        P(X) > 0 \iff X \in ]-\infty\,,\, 4[ \cup ]9\,,\, +\infty[.
        \]
        \item \textbf{Prise en compte de la contrainte $X = x^2 \ge 0$ :}
        La condition devient :
        \[
        0 \le X < 4 \quad \text{ou} \quad X > 9.
        \]
        En revenant à l'inconnue réelle $x$ :
        \begin{itemize}
            \item $0 \le x^2 < 4 \iff |x| < 2 \iff -2 < x < 2 \iff x \in ]-2\,,\, 2[$.
            \item $x^2 > 9 \iff |x| > 3 \iff x < -3 \quad \text{ou} \quad x > 3 \iff x \in ]-\infty\,,\, -3[ \cup ]3\,,\, +\infty[$.
        \end{itemize}
    \end{itemize}
    L'ensemble des solutions de la seconde inéquation est la réunion de ces intervalles disjoints :
    \[
    \mathcal{S}_2 = ]-\infty\,,\, -3[ \cup ]-2\,,\, 2[ \cup ]3\,,\, +\infty[.
    \]
\end{enumerate}
\end{correctionbox}

% ----------------------------------------------------
% EXERCICE 11
% ----------------------------------------------------
\subsection*{Exercice 11 : Équation polynomiale d'ordre supérieur réductible}
Résoudre dans $\R$ l'équation :
\[
x^7 + 3x^5 - 4x^3 = 0.
\]

\begin{correctionbox}
\begin{enumerate}
    \item \textbf{Factorisation par le monôme de plus bas degré :}
    On factorise par le terme commun $x^3$ :
    \[
    x^3 (x^4 + 3x^2 - 4) = 0.
    \]
    D'après la règle du produit nul :
    \[
    x^3 = 0 \iff x = 0
    \]
    ou
    \[
    x^4 + 3x^2 - 4 = 0.
    \]

    \item \textbf{Résolution de l'équation bicarrée $x^4 + 3x^2 - 4 = 0$ :}
    Effectuons le changement de variable $X = x^2$ sous la contrainte $X \ge 0$. L'équation devient :
    \[
    X^2 + 3X - 4 = 0.
    \]
    La somme des coefficients vaut $1 + 3 + (-4) = 0$, donc $X_1 = 1$ est une racine évidente.
    Le produit des racines étant $\frac{c}{a} = -4$, la seconde racine est $X_2 = -4$.

    \item \textbf{Retour à la variable $x$ :}
    \begin{itemize}
        \item Pour $X_2 = -4$ : l'équation $x^2 = -4$ n'admet aucune solution réelle car pour tout $x \in \R$, $x^2 \ge 0 > -4$.
        \item Pour $X_1 = 1$ : l'équation $x^2 = 1$ admet deux solutions réelles :
        \[
        x = 1 \quad \text{ou} \quad x = -1.
        \]
    \end{itemize}

    \item \textbf{Conclusion :}
    En combinant la solution nulle $x = 0$ et les solutions issues de l'équation bicarrée, on obtient :
    \[
    \mathcal{S} = \{-1\,,\, 0\,,\, 1\}.
    \]
\end{enumerate}
\end{correctionbox}

% ----------------------------------------------------
% EXERCICE 12
% ----------------------------------------------------
\subsection*{Exercice 12 : Équation symétrique (ou réciproque)}
Résoudre dans $\R$ l'équation :
\[
x^2 + x - 4 + \frac{1}{x} + \frac{1}{x^2} = 0.
\]

\begin{correctionbox}
\begin{enumerate}
    \item \textbf{Ensemble de validité :}
    Les fractions rationnelles $\frac{1}{x}$ et $\frac{1}{x^2}$ imposent $x \neq 0$.
    L'ensemble de validité de l'équation est donc :
    \[
    \mathcal{D} = \R^* = \R \setminus \{0\}.
    \]

    \item \textbf{Regroupement symétrique des termes :}
    Pour tout $x \in \mathcal{D}$, réorganisons les termes en groupant les puissances opposées :
    \[
    \left(x^2 + \frac{1}{x^2}\right) + \left(x + \frac{1}{x}\right) - 4 = 0.
    \]

    \item \textbf{Changement de variable symétrique :}
    Posons la nouvelle variable :
    \[
    X = x + \frac{1}{x}.
    \]
    Élevons cette égalité au carré à l'aide de la première identité remarquable :
    \[
    X^2 = \left(x + \frac{1}{x}\right)^2 = x^2 + 2 \times x \times \frac{1}{x} + \frac{1}{x^2} = x^2 + 2 + \frac{1}{x^2}.
    \]
    On en déduit l'expression du terme quadratique en fonction de $X$ :
    \[
    x^2 + \frac{1}{x^2} = X^2 - 2.
    \]

    \item \textbf{Résolution de l'équation auxiliaire en $X$ :}
    En substituant dans l'équation de départ :
    \[
    (X^2 - 2) + X - 4 = 0 \iff X^2 + X - 6 = 0.
    \]
    Factorisons ce trinôme en remarquant que $2$ et $-3$ ont pour somme $-1$ et produit $-6$ :
    \[
    (X - 2)(X + 3) = 0 \iff X = 2 \quad \text{ou} \quad X = -3.
    \]

    \item \textbf{Résolution en $x$ pour chaque valeur de $X$ :}
    \begin{itemize}
        \item \textbf{Premier cas : $X = 2$}
        \[
        x + \frac{1}{x} = 2 \iff \frac{x^2 - 2x + 1}{x} = 0 \iff (x - 1)^2 = 0 \iff x = 1.
        \]
        La valeur $x = 1$ appartient bien à $\mathcal{D}$.
        \item \textbf{Second cas : $X = -3$}
        \[
        x + \frac{1}{x} = -3 \iff \frac{x^2 + 3x + 1}{x} = 0 \iff x^2 + 3x + 1 = 0.
        \]
        Calculons le discriminant de ce trinôme :
        \[
        \Delta = 3^2 - 4 \times 1 \times 1 = 9 - 4 = 5 > 0.
        \]
        L'équation admet deux racines réelles distinctes :
        \[
        x_1 = \frac{-3 - \sqrt{5}}{2} \quad \text{et} \quad x_2 = \frac{-3 + \sqrt{5}}{2}.
        \]
        Ces deux valeurs sont non nulles (car $\sqrt{5} \neq 3$), elles appartiennent donc à $\mathcal{D}$.
    \end{itemize}

    \item \textbf{Conclusion :}
    L'équation admet trois solutions réelles distinctes :
    \[
    \mathcal{S} = \left\{ \frac{-3 - \sqrt{5}}{2}\,,\, \frac{-3 + \sqrt{5}}{2}\,,\, 1 \right\}.
    \]
\end{enumerate}
\end{correctionbox}

% ----------------------------------------------------
% EXERCICE 13
% ----------------------------------------------------
\subsection*{Exercice 13 : Équation à paramètre et discussion du discriminant}
Discuter selon les valeurs du paramètre réel $m$ le nombre et la nature des solutions réelles de l'équation :
\[
x^2 - (m+1)x + 4 = 0.
\]

\begin{correctionbox}
\begin{enumerate}
    \item \textbf{Nature de l'équation :}
    Pour toute valeur du paramètre $m \in \R$, le coefficient dominant devant $x^2$ vaut $1 \neq 0$. L'équation est donc toujours une équation polynomiale du second degré en $x$.

    \item \textbf{Calcul et factorisation du discriminant :}
    Le discriminant de l'équation est :
    \[
    \Delta(m) = \bigl(-(m+1)\bigr)^2 - 4 \times 1 \times 4 = (m+1)^2 - 16 = (m+1)^2 - 4^2.
    \]
    En utilisant l'identité remarquable $A^2 - B^2 = (A-B)(A+B)$ :
    \[
    \Delta(m) = (m + 1 - 4)(m + 1 + 4) = (m - 3)(m + 5).
    \]

    \item \textbf{Étude du signe de $\Delta(m)$ :}
    $\Delta(m)$ est un trinôme du second degré en $m$ dont les racines sont $-5$ et $3$. Le coefficient de $m^2$ étant $1 > 0$, $\Delta(m)$ est strictement négatif entre les racines et strictement positif à l'extérieur.

    \item \textbf{Discussion complète :}
    \begin{itemize}
        \item \textbf{Si $m \in ]-\infty\,,\, -5[ \cup ]3\,,\, +\infty[$ :}
        $\Delta(m) > 0$. L'équation admet \textbf{deux solutions réelles distinctes} :
        \[
        x_1 = \frac{m + 1 - \sqrt{(m-3)(m+5)}}{2} \quad \text{et} \quad x_2 = \frac{m + 1 + \sqrt{(m-3)(m+5)}}{2}.
        \]
        \item \textbf{Si $m = -5$ :}
        $\Delta(-5) = 0$. L'équation admet \textbf{une unique racine réelle double} :
        \[
        x_0 = \frac{m + 1}{2} = \frac{-5 + 1}{2} = -2.
        \]
        \item \textbf{Si $m = 3$ :}
        $\Delta(3) = 0$. L'équation admet \textbf{une unique racine réelle double} :
        \[
        x_0 = \frac{m + 1}{2} = \frac{3 + 1}{2} = 2.
        \]
        \item \textbf{Si $m \in ]-5\,,\, 3[$ :}
        $\Delta(m) < 0$. L'équation n'admet \textbf{aucune solution réelle} :
        \[
        \mathcal{S} = \emptyset.
        \]
    \end{itemize}
\end{enumerate}
\end{correctionbox}

% ----------------------------------------------------
% EXERCICE 14
% ----------------------------------------------------
\subsection*{Exercice 14 : Équations paramétriques et dégénérescence de degré}
Discuter selon les valeurs du paramètre réel $m$ le nombre de solutions réelles des équations suivantes :
\begin{enumerate}
    \item $(m-1)x^2 + 2x - 5 = 0$
    \item $3x^2 - mx + 4 = 0$
\end{enumerate}

\begin{correctionbox}
\begin{enumerate}
    \item \textbf{Première équation :} $(m-1)x^2 + 2x - 5 = 0$.
    Le coefficient devant le terme de plus haut degré dépend de $m$. Il est impératif de distinguer le cas où ce coefficient s'annule :
    \begin{itemize}
        \item \textbf{Cas 1 : $m = 1$ (dégénérescence du degré).}
        Le coefficient dominant est nul. L'équation n'est plus du second degré mais affine du premier degré :
        \[
        0 \cdot x^2 + 2x - 5 = 0 \iff 2x = 5 \iff x = \frac{5}{2}.
        \]
        L'équation admet alors \textbf{une unique solution réelle} : $\mathcal{S} = \left\{\frac{5}{2}\right\}$.

        \item \textbf{Cas 2 : $m \neq 1$ (équation du second degré).}
        Le coefficient dominant $m - 1$ est non nul. Calculons le discriminant :
        \[
        \Delta(m) = 2^2 - 4(m-1)(-5) = 4 + 20(m-1) = 4 + 20m - 20 = 20m - 16 = 4(5m - 4).
        \]
        Étudions le signe de $\Delta(m)$ :
        \begin{itemize}
            \item Si $5m - 4 < 0 \iff m < \frac{4}{5}$ : $\Delta(m) < 0$, l'équation n'admet \textbf{aucune solution réelle} ($\mathcal{S} = \emptyset$).
            \item Si $5m - 4 = 0 \iff m = \frac{4}{5}$ : $\Delta(m) = 0$, l'équation admet \textbf{une unique racine réelle double} :
            \[
            x_0 = \frac{-2}{2(m-1)} = \frac{-1}{\frac{4}{5} - 1} = \frac{-1}{-\frac{1}{5}} = 5.
            \]
            \item Si $m > \frac{4}{5}$ avec $m \neq 1$ : $\Delta(m) > 0$, l'équation admet \textbf{deux solutions réelles distinctes} :
            \[
            x_1 = \frac{-2 - \sqrt{4(5m-4)}}{2(m-1)} = \frac{-1 - \sqrt{5m-4}}{m-1},
            \]
            \[
            x_2 = \frac{-2 + \sqrt{4(5m-4)}}{2(m-1)} = \frac{-1 + \sqrt{5m-4}}{m-1}.
            \]
        \end{itemize}
    \end{itemize}
    \textbf{Synthèse pour l'équation (1) :}
    \begin{itemize}
        \item Si $m < \frac{4}{5}$ : $0$ solution réelle ;
        \item Si $m = \frac{4}{5}$ ou $m = 1$ : exactement $1$ solution réelle ;
        \item Si $m \in ]\frac{4}{5}\,,\, 1[ \cup ]1\,,\, +\infty[$ : exactement $2$ solutions réelles distinctes.
    \end{itemize}

    \item \textbf{Seconde équation :} $3x^2 - mx + 4 = 0$.
    Le coefficient dominant vaut $3 \neq 0$ : c'est une équation du second degré pour tout $m \in \R$.
    Calculons le discriminant :
    \[
    \Delta(m) = (-m)^2 - 4 \times 3 \times 4 = m^2 - 48 = (m - 4\sqrt{3})(m + 4\sqrt{3}).
    \]
    \begin{itemize}
        \item Si $m \in ]-\infty\,,\, -4\sqrt{3}[ \cup ]4\sqrt{3}\,,\, +\infty[$ : $\Delta(m) > 0$, l'équation admet \textbf{deux solutions réelles distinctes} :
        \[
        x_1 = \frac{m - \sqrt{m^2 - 48}}{6} \quad \text{et} \quad x_2 = \frac{m + \sqrt{m^2 - 48}}{6}.
        \]
        \item Si $m = -4\sqrt{3}$ ou $m = 4\sqrt{3}$ : $\Delta(m) = 0$, l'équation admet \textbf{une unique racine réelle double} :
        \[
        x_0 = \frac{m}{6} = \pm \frac{4\sqrt{3}}{6} = \pm \frac{2\sqrt{3}}{3}.
        \]
        \item Si $m \in ]-4\sqrt{3}\,,\, 4\sqrt{3}[$ : $\Delta(m) < 0$, l'équation n'admet \textbf{aucune solution réelle} ($\mathcal{S} = \emptyset$).
    \end{itemize}
\end{enumerate}
\end{correctionbox}

% ----------------------------------------------------
% EXERCICE 15
% ----------------------------------------------------
\subsection*{Exercice 15 : Équations à paramètres littéraux multiples}
Résoudre et discuter dans $\R$, en fonction des paramètres réels $a$ et $b$, les équations suivantes :
\begin{enumerate}
    \item $bx^2 - ax + bx - a = 0$
    \item $(ax)^2 - 3abx + 3b^2 = 0$
\end{enumerate}

\begin{correctionbox}
\begin{enumerate}
    \item \textbf{Première équation :} $bx^2 - ax + bx - a = 0$.
    Regroupons les termes par factorisation directe :
    \[
    (bx^2 + bx) - (ax + a) = bx(x + 1) - a(x + 1) = (bx - a)(x + 1) = 0.
    \]
    Discutons selon les valeurs de $b$ :
    \begin{itemize}
        \item \textbf{Si $b = 0$ :} L'équation devient $-a(x+1) = 0$.
        \begin{itemize}
            \item Si $a = 0$ : l'équation s'écrit $0 = 0$, elle est vérifiée pour tout $x \in \R$, donc $\mathcal{S} = \R$.
            \item Si $a \neq 0$ : l'équation équivaut à $x + 1 = 0$, d'où $\mathcal{S} = \{-1\}$.
        \end{itemize}
        \item \textbf{Si $b \neq 0$ :} L'équation est du second degré. Par le produit nul :
        \[
        bx - a = 0 \iff x = \frac{a}{b} \quad \text{ou} \quad x + 1 = 0 \iff x = -1.
        \]
        \begin{itemize}
            \item Si $\frac{a}{b} = -1 \iff a = -b$ : les deux racines coïncident, l'équation admet une unique racine double $x_0 = -1$, donc $\mathcal{S} = \{-1\}$.
            \item Si $a \neq -b$ : l'équation admet deux racines distinctes, donc $\mathcal{S} = \left\{-1\,,\, \frac{a}{b}\right\}$.
        \end{itemize}
    \end{itemize}
    \textbf{Remarque par le discriminant :} Sous forme développée $bx^2 + (b-a)x - a = 0$, le discriminant est :
    \[
    \Delta = (b-a)^2 - 4b(-a) = (b-a)^2 + 4ab = b^2 - 2ab + a^2 + 4ab = (a+b)^2 \ge 0.
    \]
    Comme $\Delta \ge 0$, l'équation admet toujours au moins une solution réelle pour $b \neq 0$.

    \item \textbf{Seconde équation :} $(ax)^2 - 3abx + 3b^2 = 0 \iff a^2 x^2 - 3abx + 3b^2 = 0$.
    \begin{itemize}
        \item \textbf{Si $a = 0$ :} L'équation se réduit à :
        \[
        3b^2 = 0 \iff b = 0.
        \]
        \begin{itemize}
            \item Si $b = 0$ (donc $a = 0$ et $b = 0$) : $0 = 0$ est toujours vrai, donc $\mathcal{S} = \R$.
            \item Si $b \neq 0$ (avec $a = 0$) : l'équation $3b^2 = 0$ n'a aucune solution, donc $\mathcal{S} = \emptyset$.
        \end{itemize}
        \item \textbf{Si $a \neq 0$ :} Le coefficient dominant est $a^2 > 0$ ; l'équation est du second degré.
        Calculons son discriminant :
        \[
        \Delta = (-3ab)^2 - 4(a^2)(3b^2) = 9a^2 b^2 - 12a^2 b^2 = -3a^2 b^2 = -3(ab)^2.
        \]
        Puisqu'un carré est toujours positif ou nul, $(ab)^2 \ge 0$, donc :
        \[
        \Delta = -3(ab)^2 \le 0.
        \]
        \begin{itemize}
            \item Si $b = 0$ : alors $\Delta = 0$. L'équation devient $a^2 x^2 = 0 \iff x = 0$ (racine double unique), donc $\mathcal{S} = \{0\}$.
            \item Si $b \neq 0$ : comme $a \neq 0$, on a $ab \neq 0$, donc $\Delta = -3(ab)^2 < 0$.
            L'équation n'admet aucune solution réelle : $\mathcal{S} = \emptyset$.
        \end{itemize}
    \end{itemize}
\end{enumerate}
\end{correctionbox}

% ----------------------------------------------------
% EXERCICE 16
% ----------------------------------------------------
\subsection*{Exercice 16 : Existence garantie de racines réelles pour une équation littérale}
Soient $a, b, c$ trois nombres réels fixés.
\begin{enumerate}
    \item Démontrer que l'équation d'inconnue $x$ :
    \[
    (x-a)(x-b) = c^2
    \]
    admet au moins une solution réelle pour tout triplet $(a, b, c) \in \R^3$.
    \item À quelle condition nécessaire et suffisante sur $a, b, c$ l'équation admet-elle une unique solution réelle ?
\end{enumerate}

\begin{correctionbox}
\begin{enumerate}
    \item Développons le membre de gauche et rassemblons tous les termes :
    \[
    (x-a)(x-b) = c^2 \iff x^2 - (a+b)x + ab - c^2 = 0.
    \]
    Il s'agit d'une équation polynomiale du second degré en $x$, dont le coefficient dominant est $1 \neq 0$.
    Calculons son discriminant :
    \[
    \Delta = \bigl(-(a+b)\bigr)^2 - 4 \times 1 \times (ab - c^2) = (a+b)^2 - 4ab + 4c^2.
    \]
    Développons le premier bloc :
    \[
    (a+b)^2 - 4ab = a^2 + 2ab + b^2 - 4ab = a^2 - 2ab + b^2 = (a-b)^2.
    \]
    Le discriminant s'écrit donc sous la forme remarquable d'une somme de deux carrés :
    \[
    \Delta = (a-b)^2 + 4c^2.
    \]
    Or, pour tous réels $a, b, c$, les carrés $(a-b)^2 \ge 0$ et $c^2 \ge 0$ sont positifs ou nuls.
    Par conséquent :
    \[
    \forall (a, b, c) \in \R^3, \quad \Delta = (a-b)^2 + 4c^2 \ge 0.
    \]
    Le discriminant étant positif ou nul pour tout choix de paramètres réels, l'équation admet toujours au moins une racine réelle.

    \item Une équation du second degré admet une unique solution réelle (racine double) si et seulement si son discriminant est strictement nul :
    \[
    \Delta = 0 \iff (a-b)^2 + 4c^2 = 0.
    \]
    Une somme de deux nombres réels positifs ou nuls est nulle si et seulement si chacun d'eux est nul :
    \[
    \begin{cases}
    (a-b)^2 = 0 \\
    4c^2 = 0
    \end{cases}
    \iff
    \begin{cases}
    a - b = 0 \\
    c = 0
    \end{cases}
    \iff
    \begin{cases}
    a = b \\
    c = 0.
    \end{cases}
    \]
    Lorsque cette condition est remplie, l'unique racine double est :
    \[
    x_0 = \frac{a+b}{2} = a.
    \]
    \textbf{Conclusion :} L'équation admet une unique solution réelle si et seulement si $a = b$ et $c = 0$.
\end{enumerate}
\end{correctionbox}

% ----------------------------------------------------
% EXERCICE 17
% ----------------------------------------------------
\subsection*{Exercice 17 : Radicaux imbriqués et trinôme (I)}
\begin{enumerate}
    \item Développer et simplifier le nombre $(4\sqrt{3} - \sqrt{6})^2$.
    \item En déduire la résolution dans $\R$ de l'équation :
    \[
    4x^2 - (\sqrt{6} + 4\sqrt{3})x + \sqrt{18} = 0.
    \]
\end{enumerate}

\begin{correctionbox}
\begin{enumerate}
    \item Développons le carré $(4\sqrt{3} - \sqrt{6})^2$ à l'aide de l'identité remarquable $(u-v)^2 = u^2 - 2uv + v^2$ :
    \begin{align*}
    (4\sqrt{3} - \sqrt{6})^2 &= (4\sqrt{3})^2 - 2(4\sqrt{3})(\sqrt{6}) + (\sqrt{6})^2 \\
    &= 16 \times 3 - 8\sqrt{18} + 6 \\
    &= 48 - 8\sqrt{9 \times 2} + 6 \\
    &= 54 - 8 \times 3\sqrt{2} \\
    &= 54 - 24\sqrt{2}.
    \end{align*}

    \item Résolvons l'équation du second degré $4x^2 - (\sqrt{6} + 4\sqrt{3})x + \sqrt{18} = 0$.
    Ici, $a = 4$, $b = -(\sqrt{6} + 4\sqrt{3})$ et $c = \sqrt{18} = 3\sqrt{2}$.
    Calculons le discriminant :
    \[
    \Delta = \bigl(-(\sqrt{6} + 4\sqrt{3})\bigr)^2 - 4 \times 4 \times \sqrt{18} = (\sqrt{6} + 4\sqrt{3})^2 - 16\sqrt{18}.
    \]
    Développons le premier terme :
    \[
    (\sqrt{6} + 4\sqrt{3})^2 = 6 + 2 \times 4\sqrt{18} + 48 = 54 + 8\sqrt{18} = 54 + 24\sqrt{2}.
    \]
    Comme $16\sqrt{18} = 16 \times 3\sqrt{2} = 48\sqrt{2}$, il vient :
    \[
    \Delta = (54 + 24\sqrt{2}) - 48\sqrt{2} = 54 - 24\sqrt{2}.
    \]
    D'après le calcul de la question 1, on identifie immédiatement :
    \[
    \Delta = (4\sqrt{3} - \sqrt{6})^2.
    \]
    Comme $4\sqrt{3} = \sqrt{48} > \sqrt{6}$, la quantité $4\sqrt{3} - \sqrt{6}$ est strictement positive. Par conséquent :
    \[
    \sqrt{\Delta} = 4\sqrt{3} - \sqrt{6}.
    \]
    Les deux racines réelles sont données par :
    \[
    x_1 = \frac{(\sqrt{6} + 4\sqrt{3}) - (4\sqrt{3} - \sqrt{6})}{2 \times 4} = \frac{2\sqrt{6}}{8} = \frac{\sqrt{6}}{4},
    \]
    \[
    x_2 = \frac{(\sqrt{6} + 4\sqrt{3}) + (4\sqrt{3} - \sqrt{6})}{2 \times 4} = \frac{8\sqrt{3}}{8} = \sqrt{3}.
    \]

    \item \textbf{Conclusion :}
    L'ensemble des solutions de l'équation est :
    \[
    \mathcal{S} = \left\{ \frac{\sqrt{6}}{4}\,,\, \sqrt{3} \right\}.
    \]
\end{enumerate}
\end{correctionbox}

% ----------------------------------------------------
% EXERCICE 18
% ----------------------------------------------------
\subsection*{Exercice 18 : Radicaux imbriqués et trinôme (II)}
\begin{enumerate}
    \item Développer et simplifier le nombre $(3 - \sqrt{2})^2$.
    \item Résoudre dans $\R$ l'équation :
    \[
    x^2 + (1 + \sqrt{2})x + 2(\sqrt{2} - 1) = 0.
    \]
\end{enumerate}

\begin{correctionbox}
\begin{enumerate}
    \item Développons le carré $(3 - \sqrt{2})^2$ :
    \[
    (3 - \sqrt{2})^2 = 3^2 - 2 \times 3 \times \sqrt{2} + (\sqrt{2})^2 = 9 - 6\sqrt{2} + 2 = 11 - 6\sqrt{2}.
    \]

    \item Résolvons l'équation $x^2 + (1 + \sqrt{2})x + 2(\sqrt{2} - 1) = 0$.
    Ici $a = 1$, $b = 1 + \sqrt{2}$ et $c = 2(\sqrt{2} - 1)$.
    Calculons son discriminant :
    \[
    \Delta = (1 + \sqrt{2})^2 - 4 \times 1 \times 2(\sqrt{2} - 1) = (1 + \sqrt{2})^2 - 8(\sqrt{2} - 1).
    \]
    Développons les deux expressions :
    \[
    (1 + \sqrt{2})^2 = 1 + 2\sqrt{2} + 2 = 3 + 2\sqrt{2},
    \]
    \[
    -8(\sqrt{2} - 1) = -8\sqrt{2} + 8.
    \]
    En additionnant ces résultats :
    \[
    \Delta = (3 + 2\sqrt{2}) + (-8\sqrt{2} + 8) = 11 - 6\sqrt{2}.
    \]
    D'après la question 1, on reconnaît :
    \[
    \Delta = (3 - \sqrt{2})^2.
    \]
    Comme $3 > \sqrt{2} \approx 1{,}414$, on a $3 - \sqrt{2} > 0$, d'où :
    \[
    \sqrt{\Delta} = 3 - \sqrt{2}.
    \]
    Les deux racines réelles de l'équation sont :
    \[
    x_1 = \frac{-(1 + \sqrt{2}) - (3 - \sqrt{2})}{2} = \frac{-1 - \sqrt{2} - 3 + \sqrt{2}}{2} = \frac{-4}{2} = -2,
    \]
    \[
    x_2 = \frac{-(1 + \sqrt{2}) + (3 - \sqrt{2})}{2} = \frac{-1 - \sqrt{2} + 3 - \sqrt{2}}{2} = \frac{2 - 2\sqrt{2}}{2} = 1 - \sqrt{2}.
    \]

    \item \textbf{Conclusion :}
    L'ensemble des solutions est :
    \[
    \mathcal{S} = \{-2\,,\, 1 - \sqrt{2}\}.
    \]
\end{enumerate}
\end{correctionbox}

% ----------------------------------------------------
% EXERCICE 19
% ----------------------------------------------------
\subsection*{Exercice 19 : Le rectangle d'or et la proportion divine}
Un rectangle de longueur $L$ et de largeur $\ell$ (avec $L > \ell > 0$) est appelé « rectangle d'or » lorsque, après lui avoir retiré un carré de côté $\ell$, le rectangle résiduel (de dimensions $\ell$ et $L - \ell$) est semblable au rectangle initial, c'est-à-dire que le rapport de sa longueur sur sa largeur est conservé.
On pose $\Phi = \frac{L}{\ell}$.
\begin{enumerate}
    \item Montrer que $\Phi$ est l'unique racine strictement positive de l'équation $x^2 - x - 1 = 0$.
    \item En déduire la valeur exacte de $\Phi$ (le « nombre d'or »).
    \item Établir les relations algébriques remarquables suivantes :
    \[
    \frac{1}{\Phi} = \Phi - 1, \quad \Phi^2 = \Phi + 1, \quad \Phi^3 = 2\Phi + 1 \quad \text{et} \quad \Phi^4 = 3\Phi + 2.
    \]
\end{enumerate}

\begin{correctionbox}
\begin{enumerate}
    \item \textbf{Mise en équation géométrique :}
    Dans le rectangle initial, le rapport longueur sur largeur est $\frac{L}{\ell} = \Phi > 1$.
    Le rectangle restant a pour côtés $\ell$ et $L - \ell$.
    Sa plus grande dimension étant $\ell$ et sa plus petite étant $L - \ell$, la condition de proportionnalité (similitude) s'énonce :
    \[
    \frac{L}{\ell} = \frac{\ell}{L - \ell}.
    \]
    Divisons par $\ell$ le numérateur et le dénominateur du membre de droite :
    \[
    \frac{\ell}{L - \ell} = \frac{1}{\frac{L}{\ell} - 1} = \frac{1}{\Phi - 1}.
    \]
    La relation devient ainsi :
    \[
    \Phi = \frac{1}{\Phi - 1}.
    \]
    Puisque $\Phi > 1$, on peut multiplier par $\Phi - 1$ sans introduire d'équivoque :
    \[
    \Phi(\Phi - 1) = 1 \iff \Phi^2 - \Phi - 1 = 0.
    \]
    Le réel $\Phi$ est donc racine de l'équation du second degré $x^2 - x - 1 = 0$.

    \item \textbf{Calcul de la valeur exacte du nombre d'or :}
    Résolvons l'équation $x^2 - x - 1 = 0$.
    Son discriminant est :
    \[
    \Delta = (-1)^2 - 4 \times 1 \times (-1) = 1 + 4 = 5 > 0.
    \]
    L'équation admet deux racines réelles distinctes :
    \[
    x_1 = \frac{1 - \sqrt{5}}{2} \quad \text{et} \quad x_2 = \frac{1 + \sqrt{5}}{2}.
    \]
    Comme $\sqrt{5} \approx 2{,}236 > 1$, la racine $x_1$ est strictement négative ($x_1 \approx -0{,}618$).
    Puisque $\Phi = \frac{L}{\ell}$ est le rapport de deux grandeurs géométriques strictement positives, on a impérativement $\Phi > 0$.
    On en déduit que le nombre d'or est :
    \[
    \Phi = \frac{1 + \sqrt{5}}{2} \approx 1{,}6180339887\dots
    \]

    \item \textbf{Relations algébriques et puissances de $\Phi$ :}
    \begin{itemize}
        \item De l'égalité $\Phi^2 - \Phi - 1 = 0$, en divisant par $\Phi > 0$ :
        \[
        \Phi - 1 - \frac{1}{\Phi} = 0 \iff \frac{1}{\Phi} = \Phi - 1.
        \]
        \item En isolant $\Phi^2$ dans l'équation :
        \[
        \Phi^2 = \Phi + 1.
        \]
        \item En multipliant cette relation par $\Phi$ :
        \[
        \Phi^3 = \Phi \times \Phi^2 = \Phi(\Phi + 1) = \Phi^2 + \Phi = (\Phi + 1) + \Phi = 2\Phi + 1.
        \]
        \item En multipliant à nouveau par $\Phi$ :
        \[
        \Phi^4 = \Phi(2\Phi + 1) = 2\Phi^2 + \Phi = 2(\Phi + 1) + \Phi = 3\Phi + 2.
        \]
    \end{itemize}
    \textbf{Remarque culturelle :} On reconnaît les nombres de Fibonacci $(F_n)_{n \in \N} = (1, 1, 2, 3, 5, 8, \dots)$ tels que pour tout $n \ge 1$ :
    \[
    \Phi^n = F_n \Phi + F_{n-1}.
    \]
\end{enumerate}
\end{correctionbox}

% ----------------------------------------------------
% EXERCICE 20
% ----------------------------------------------------
\subsection*{Exercice 20 : Radical imbriqué infini et point fixe}
On considère l'expression formelle définie par une infinité de radicaux imbriqués :
\[
A = \sqrt{1 + \sqrt{1 + \sqrt{1 + \sqrt{1 + \cdots}}}}
\]
Afin de donner une assise analytique rigoureuse à cet objet, on introduit la suite réelle $(u_n)_{n \in \N}$ définie par :
\[
u_0 = 1 \quad \text{et} \quad \forall n \in \N, \quad u_{n+1} = \sqrt{1 + u_n}.
\]
\begin{enumerate}
    \item Démontrer par récurrence que pour tout entier naturel $n$, on a : $1 \le u_n \le 2$.
    \item Démontrer que la suite $(u_n)_{n \in \N}$ est strictement croissante.
    \item Justifier que la suite $(u_n)_{n \in \N}$ converge vers une limite finie $\ell \in [1\,,\, 2]$.
    \item Déterminer la valeur exacte de $\ell$ et faire le lien avec le nombre d'or $\Phi$.
\end{enumerate}

\begin{correctionbox}
\begin{enumerate}
    \item \textbf{Démonstration par récurrence du bornage :}
    Pour tout $n \in \N$, notons la propriété $\mathcal{P}(n) : 1 \le u_n \le 2$.
    \begin{itemize}
        \item \textbf{Initialisation :} Pour $n = 0$, $u_0 = 1$. L'encadrement $1 \le 1 \le 2$ est vérifié, donc $\mathcal{P}(0)$ est vraie.
        \item \textbf{Hérédité :} Soit $n \in \N$. Supposons $\mathcal{P}(n)$ vraie, c'est-à-dire $1 \le u_n \le 2$.
        Ajoutons $1$ aux trois membres de cette inégalité :
        \[
        2 \le 1 + u_n \le 3.
        \]
        La fonction racine carrée $x \mapsto \sqrt{x}$ étant strictement croissante sur $[0\,,\, +\infty[$ :
        \[
        \sqrt{2} \le \sqrt{1 + u_n} \le \sqrt{3} \iff \sqrt{2} \le u_{n+1} \le \sqrt{3}.
        \]
        Or, on a clairement $1 < \sqrt{2}$ et $\sqrt{3} < \sqrt{4} = 2$.
        On en déduit donc l'encadrement :
        \[
        1 \le u_{n+1} \le 2,
        \]
        ce qui prouve la véracité de $\mathcal{P}(n+1)$.
    \end{itemize}
    Par le principe de récurrence simple, on conclut que :
    \[
    \forall n \in \N, \quad 1 \le u_n \le 2.
    \]

    \item \textbf{Stricte croissance de la suite $(u_n)_{n \in \N}$ :}
    Démontrons par récurrence sur $n \in \N$ la propriété $\mathcal{Q}(n) : u_n < u_{n+1}$.
    \begin{itemize}
        \item \textbf{Initialisation :} Pour $n = 0$, on a $u_0 = 1$ et $u_1 = \sqrt{1 + 1} = \sqrt{2} \approx 1{,}414$.
        Comme $1 < \sqrt{2}$, la relation $u_0 < u_1$ est vérifiée. Donc $\mathcal{Q}(0)$ est vraie.
        \item \textbf{Hérédité :} Soit $n \in \N$. Supposons que $\mathcal{Q}(n)$ soit vraie, c'est-à-dire $u_n < u_{n+1}$.
        En ajoutant $1$ :
        \[
        1 + u_n < 1 + u_{n+1}.
        \]
        Par stricte croissance de la fonction racine carrée sur $[0, +\infty[$ :
        \[
        \sqrt{1 + u_n} < \sqrt{1 + u_{n+1}} \iff u_{n+1} < u_{n+2}.
        \]
        La propriété $\mathcal{Q}(n+1)$ est ainsi établie.
    \end{itemize}
    Par récurrence, on conclut que pour tout $n \in \N$, $u_n < u_{n+1}$.
    La suite $(u_n)_{n \in \N}$ est donc strictement croissante.

    \item \textbf{Convergence de la suite :}
    La suite réelle $(u_n)_{n \in \N}$ est :
    \begin{itemize}
        \item strictement croissante ;
        \item majorée par $2$.
    \end{itemize}
    D'après le \textbf{théorème de convergence monotone}, toute suite croissante et majorée converge dans $\R$.
    Ainsi, $(u_n)_{n \in \N}$ converge vers une limite finie $\ell \in \R$.
    De plus, comme pour tout $n \in \N$, $1 \le u_n \le 2$, le passage à la limite dans les inégalités larges garantit que :
    \[
    1 \le \ell \le 2.
    \]

    \item \textbf{Détermination de la limite $\ell$ et lien avec le nombre d'or :}
    La fonction $f : x \mapsto \sqrt{1 + x}$ est continue sur son domaine de définition $[-1\,,\, +\infty[$, donc en particulier sur l'intervalle $[1\,,\, 2]$.
    En faisant tendre $n$ vers $+\infty$ dans la relation de récurrence $u_{n+1} = f(u_n)$, on obtient par continuité de $f$ l'équation de point fixe :
    \[
    \ell = f(\ell) \iff \ell = \sqrt{1 + \ell}.
    \]
    Puisque $\ell \ge 1 > 0$, les deux membres sont strictement positifs. On peut élever au carré par équivalence :
    \[
    \ell^2 = 1 + \ell \iff \ell^2 - \ell - 1 = 0.
    \]
    On retrouve l'équation caractéristique du nombre d'or étudiée à l'exercice 19.
    Les deux solutions de cette équation sont :
    \[
    \ell_1 = \frac{1 - \sqrt{5}}{2} < 0 \quad \text{et} \quad \ell_2 = \frac{1 + \sqrt{5}}{2} > 0.
    \]
    Comme la limite $\ell$ doit vérifier $\ell \ge 1 > 0$, la racine négative $\ell_1$ est écartée.
    L'unique limite possible est donc :
    \[
    \ell = \frac{1 + \sqrt{5}}{2} = \Phi.
    \]

    \item \textbf{Conclusion :}
    La suite converge vers le nombre d'or $\Phi$. Par conséquent, la quantité formelle $A$ s'interprète rigoureusement comme la limite de la suite $(u_n)$, et vaut exactement :
    \[
    A = \lim_{n \to +\infty} u_n = \frac{1 + \sqrt{5}}{2} = \Phi \approx 1{,}6180339887\dots
    \]
\end{enumerate}
\end{correctionbox}


\section{Systèmes d'équations et modélisation}

% ----------------------------------------------------
% EXERCICE 21
% ----------------------------------------------------
\subsection*{Exercice 21}
Une parabole $\mathcal{C}_f$ représentant une fonction polynôme du second degré $f$ passe par les points $A(4\,;\,0)$, $B(7\,;\,0)$ et $C(0\,;\,-14)$. Déterminer l'expression de $f(x)$ sous forme développée.

\begin{correctionbox}
\textbf{1. Exploitation des points d'intersection avec l'axe des abscisses :}\\
La parabole $\mathcal{C}_f$ coupe l'axe des abscisses aux points $A(4\,;\,0)$ et $B(7\,;\,0)$. Les réels $4$ et $7$ sont donc les deux racines distinctes du trinôme du second degré $f$.

Par conséquent, d'après le théorème de factorisation des fonctions polynômes du second degré, $f(x)$ admet une factorisation de la forme :
\[
f(x) = a(x - 4)(x - 7), \quad \text{où } a \in \mathbb{R}^*.
\]

\textbf{2. Détermination du coefficient dominant $a$ :}\\
Le point $C(0\,;\,-14)$ appartient à la courbe $\mathcal{C}_f$, donc ses coordonnées vérifient l'équation de la parabole, soit $f(0) = -14$.
En substituant $x = 0$ dans l'expression factorisée :
\[
a(0 - 4)(0 - 7) = -14 \iff a \times (-4) \times (-7) = -14 \iff 28a = -14 \iff a = -\frac{14}{28} = -\frac{1}{2}.
\]

\textbf{3. Obtention de la forme développée :}\\
On développe le produit :
\[
f(x) = -\frac{1}{2}(x - 4)(x - 7) = -\frac{1}{2}(x^2 - 11x + 28) = -\frac{1}{2}x^2 + \frac{11}{2}x - 14.
\]

\textbf{Conclusion et vérification :}\\
L'expression développée cherchée est :
\[
\boxed{f(x) = -\frac{1}{2}x^2 + \frac{11}{2}x - 14}
\]
\textit{Vérification rapide :}
\begin{itemize}[leftmargin=*,label=\textbullet]
    \item $f(4) = -\frac{1}{2}(16) + \frac{11}{2}(4) - 14 = -8 + 22 - 14 = 0$ ;
    \item $f(7) = -\frac{1}{2}(49) + \frac{77}{2} - 14 = \frac{28}{2} - 14 = 14 - 14 = 0$ ;
    \item $f(0) = -14$.
\end{itemize}
Toutes les conditions sont rigoureusement satisfaites.
\end{correctionbox}

% ----------------------------------------------------
% EXERCICE 22
% ----------------------------------------------------
\subsection*{Exercice 22}
Résoudre dans $\mathbb{R}^2$ le système suivant :
\[
\begin{cases}
x^2 + y^2 = \dfrac{82}{9} \\[2mm]
xy = 1
\end{cases}
\]

\begin{correctionbox}
Soit $(x, y) \in \mathbb{R}^2$ un couple solution. Puisque $xy = 1 \ne 0$, les réels $x$ et $y$ sont nécessairement non nuls et ont le même signe.

\textbf{Méthode 1 : Utilisation des identités remarquables symétriques (très élégante)}\\
Exprimons $(x+y)^2$ et $(x-y)^2$ à l'aide des deux équations du système :
\[
\begin{aligned}
(x+y)^2 &= x^2 + y^2 + 2xy = \frac{82}{9} + 2 \times 1 = \frac{82 + 18}{9} = \frac{100}{9} = \left(\frac{10}{3}\right)^2, \\[2mm]
(x-y)^2 &= x^2 + y^2 - 2xy = \frac{82}{9} - 2 \times 1 = \frac{82 - 18}{9} = \frac{64}{9} = \left(\frac{8}{3}\right)^2.
\end{aligned}
\]
On en déduit :
\[
x+y = \pm \frac{10}{3} \quad \text{et} \quad x-y = \pm \frac{8}{3}.
\]
Les réels $x$ et $y$ ont pour produit $P = xy = 1$ et pour somme $S = x+y \in \left\{ -\frac{10}{3}, \frac{10}{3} \right\}$.
Ils sont donc les racines de l'équation du second degré $X^2 - SX + P = 0$.

\begin{itemize}[leftmargin=*,label=\textbullet]
    \item \textbf{Premier cas : $S = \frac{10}{3}$ et $P = 1$}\\
    L'équation s'écrit :
    \[
    X^2 - \frac{10}{3}X + 1 = 0 \iff 3X^2 - 10X + 3 = 0.
    \]
    Discriminant : $\Delta = (-10)^2 - 4 \times 3 \times 3 = 100 - 36 = 64 = 8^2 > 0$.\\
    Les racines sont :
    \[
    X_1 = \frac{10 - 8}{6} = \frac{2}{6} = \frac{1}{3} \quad \text{et} \quad X_2 = \frac{10 + 8}{6} = \frac{18}{6} = 3.
    \]
    On obtient deux couples de solutions : $\left(3\,;\,\frac{1}{3}\right)$ et $\left(\frac{1}{3}\,;\,3\right)$.

    \item \textbf{Deuxième cas : $S = -\frac{10}{3}$ et $P = 1$}\\
    L'équation s'écrit :
    \[
    X^2 + \frac{10}{3}X + 1 = 0 \iff 3X^2 + 10X + 3 = 0.
    \]
    Les racines sont les opposées du cas précédent :
    \[
    X_3 = -3 \quad \text{et} \quad X_4 = -\frac{1}{3}.
    \]
    On obtient deux autres couples : $\left(-3\,;\,-\frac{1}{3}\right)$ et $\left(-\frac{1}{3}\,;\,-3\right)$.
\end{itemize}

\textbf{Méthode 2 : Par substitution et équation bicarrée}\\
Puisque $x \ne 0$, la seconde équation donne $y = \frac{1}{x}$. En substituant dans la première :
\[
x^2 + \frac{1}{x^2} = \frac{82}{9} \iff x^4 - \frac{82}{9}x^2 + 1 = 0 \iff 9(x^2)^2 - 82x^2 + 9 = 0.
\]
En posant $u = x^2 \ge 0$, on obtient l'équation du second degré $9u^2 - 82u + 9 = 0$.\\
Son discriminant vaut $\Delta = (-82)^2 - 4 \times 9 \times 9 = 6724 - 324 = 6400 = 80^2$.\\
Les racines sont :
\[
u_1 = \frac{82 - 80}{18} = \frac{2}{18} = \frac{1}{9} \quad \text{et} \quad u_2 = \frac{82 + 80}{18} = \frac{162}{18} = 9.
\]
D'où :
\[
x^2 = 9 \iff x \in \{-3, 3\} \quad \text{ou} \quad x^2 = \frac{1}{9} \iff x \in \left\{ -\frac{1}{3}, \frac{1}{3} \right\}.
\]
Comme $y = \frac{1}{x}$, on retrouve exactement les quatre couples.

\textbf{Conclusion :}\\
L'ensemble des solutions du système dans $\mathbb{R}^2$ est :
\[
\boxed{\mathcal{S} = \left\{ \left(3\,;\,\frac{1}{3}\right),\, \left(\frac{1}{3}\,;\,3\right),\, \left(-3\,;\,-\frac{1}{3}\right),\, \left(-\frac{1}{3}\,;\,-3\right) \right\}}
\]
\end{correctionbox}

% ----------------------------------------------------
% EXERCICE 23
% ----------------------------------------------------
\subsection*{Exercice 23}
Calculer les longueurs des côtés de l'angle droit d'un triangle rectangle dont l'aire est de $84\text{ cm}^2$ et dont l'hypoténuse a pour mesure $25\text{ cm}$.

\begin{correctionbox}
\textbf{1. Mise en équation :}\\
Soient $x$ et $y$ les longueurs (exprimées en $\text{cm}$) des deux côtés de l'angle droit. Ce sont des réels strictement positifs ($x > 0$ et $y > 0$).
\begin{itemize}[leftmargin=*,label=\textbullet]
    \item L'aire d'un triangle rectangle est donnée par $\mathcal{A} = \frac{1}{2}xy$. L'énoncé donne :
    \[
    \frac{1}{2}xy = 84 \iff xy = 168.
    \]
    \item D'après le théorème de Pythagore, la somme des carrés des côtés de l'angle droit est égale au carré de l'hypoténuse :
    \[
    x^2 + y^2 = 25^2 = 625.
    \]
\end{itemize}
Le problème équivaut donc à résoudre le système symétrique :
\[
\begin{cases}
x^2 + y^2 = 625 \\
xy = 168
\end{cases}
\quad \text{avec } x > 0, \, y > 0.
\]

\textbf{2. Détermination de la somme $S = x + y$ :}\\
En utilisant l'identité remarquable du carré d'une somme :
\[
(x + y)^2 = x^2 + y^2 + 2xy = 625 + 2 \times 168 = 625 + 336 = 961.
\]
Calculons la racine carrée de $961$ : on remarque que $30^2 = 900$ et $31^2 = (30+1)^2 = 900 + 60 + 1 = 961$.\\
Comme $x > 0$ et $y > 0$, la somme $x+y$ est strictement positive, d'où :
\[
x + y = \sqrt{961} = 31.
\]

\textbf{3. Résolution par l'équation produit-somme :}\\
Les réels $x$ et $y$ ont pour somme $S = 31$ et pour produit $P = 168$. Ils sont donc les deux solutions de l'équation polynomiale du second degré :
\[
X^2 - 31X + 168 = 0.
\]
Calculons le discriminant :
\[
\Delta = (-31)^2 - 4 \times 1 \times 168 = 961 - 672 = 289.
\]
On reconnaît $289 = 17^2$ (puisque $17^2 = 289$).\\
Puisque $\Delta > 0$, l'équation admet deux racines réelles distinctes :
\[
X_1 = \frac{31 - 17}{2} = \frac{14}{2} = 7, \qquad X_2 = \frac{31 + 17}{2} = \frac{48}{2} = 24.
\]

\textbf{Conclusion :}\\
Les côtés de l'angle droit du triangle mesurent \textbf{$7\text{ cm}$} et \textbf{$24\text{ cm}$}.\\
\textit{Vérification :} $7^2 + 24^2 = 49 + 576 = 625 = 25^2$, et $\frac{1}{2} \times 7 \times 24 = 7 \times 12 = 84\text{ cm}^2$.
\end{correctionbox}

% ----------------------------------------------------
% EXERCICE 24
% ----------------------------------------------------
\subsection*{Exercice 24}
La somme des deux chiffres d'un entier naturel $x$ est $13$. En ajoutant $34$ à leur produit, on obtient un nombre dont les chiffres sont ceux de $x$ dans l'ordre inverse. Trouver $x$.

\begin{correctionbox}
\textbf{1. Choix des inconnues et modélisation :}\\
Soit $x$ un entier naturel s'écrivant avec deux chiffres en base 10.
Notons $a$ le chiffre des dizaines et $b$ le chiffre des unités.
Par définition de l'écriture décimale :
\[
x = 10a + b, \quad \text{avec } a \in \{1, 2, \dots, 9\} \text{ et } b \in \{0, 1, \dots, 9\}.
\]
Le nombre formé par les mêmes chiffres dans l'ordre inverse s'écrit $10b + a$.

\textbf{2. Traduction des contraintes de l'énoncé :}\\
L'énoncé fournit deux conditions :
\begin{enumerate}
    \item La somme des deux chiffres est égale à $13$ :
    \[
    a + b = 13 \iff b = 13 - a.
    \]
    \item En ajoutant $34$ au produit des deux chiffres, on obtient le nombre inversé :
    \[
    ab + 34 = 10b + a.
    \]
\end{enumerate}

\textbf{3. Résolution par substitution :}\\
Remplaçons $b$ par $13 - a$ dans la seconde équation :
\[
a(13 - a) + 34 = 10(13 - a) + a.
\]
Développons et réduisons chaque membre :
\[
\begin{aligned}
13a - a^2 + 34 &= 130 - 10a + a \\
-a^2 + 13a + 34 &= 130 - 9a.
\end{aligned}
\]
Transposons tous les termes dans le membre de droite pour obtenir un coefficient dominant positif :
\[
a^2 - 9a - 13a + 130 - 34 = 0 \iff a^2 - 22a + 96 = 0.
\]

\textbf{4. Résolution de l'équation du second degré :}\\
Discriminant du trinôme :
\[
\Delta = (-22)^2 - 4 \times 1 \times 96 = 484 - 384 = 100 = 10^2 > 0.
\]
Les deux racines réelles sont :
\[
a_1 = \frac{22 - 10}{2} = \frac{12}{2} = 6, \qquad a_2 = \frac{22 + 10}{2} = \frac{32}{2} = 16.
\]
Puisque $a$ est un chiffre entre $1$ et $9$, la solution $a_2 = 16$ est à rejeter. On retient donc uniquement :
\[
a = 6.
\]
On en déduit la valeur du chiffre des unités :
\[
b = 13 - a = 13 - 6 = 7.
\]

\textbf{Conclusion et vérification :}\\
L'entier naturel cherché est :
\[
\boxed{x = 67}
\]
\textit{Vérification :}
\begin{itemize}[leftmargin=*,label=\textbullet]
    \item Somme des chiffres : $6 + 7 = 13$ (conforme) ;
    \item Produit des chiffres augmenté de $34$ : $6 \times 7 + 34 = 42 + 34 = 76$ ;
    \item $76$ est bien le nombre formé en inversant les chiffres de $67$.
\end{itemize}
\end{correctionbox}

% ----------------------------------------------------
% EXERCICE 25
% ----------------------------------------------------
\subsection*{Exercice 25}
Deux résistances en dérivation ont une résistance équivalente de $1{,}5\,\Omega$ et en série une résistance équivalente de $8\,\Omega$. Déterminer les deux résistances.

\begin{correctionbox}
\textbf{1. Rappels physiques et formulation mathématique :}\\
Soient $R_1 > 0$ et $R_2 > 0$ les valeurs (exprimées en ohms, $\Omega$) des deux résistances.
\begin{itemize}[leftmargin=*,label=\textbullet]
    \item \textbf{Association en série :} la résistance équivalente est la somme des résistances :
    \[
    R_{\text{série}} = R_1 + R_2 = 8.
    \]
    \item \textbf{Association en dérivation (parallèle) :} l'inverse de la résistance équivalente est la somme des inverses :
    \[
    \frac{1}{R_{\text{dér}}} = \frac{1}{R_1} + \frac{1}{R_2} = \frac{R_1 + R_2}{R_1 R_2} \iff R_{\text{dér}} = \frac{R_1 R_2}{R_1 + R_2} = 1{,}5.
    \]
\end{itemize}

\textbf{2. Détermination du produit $R_1 R_2$ :}\\
En remplaçant $R_1 + R_2$ par $8$ dans l'équation de dérivation :
\[
\frac{R_1 R_2}{8} = 1{,}5 \iff R_1 R_2 = 1{,}5 \times 8 = 12.
\]

\textbf{3. Résolution du système somme-produit :}\\
Les réels strictement positifs $R_1$ et $R_2$ ont pour somme $S = 8$ et pour produit $P = 12$. Ils sont donc les racines de l'équation polynomiale du second degré :
\[
X^2 - 8X + 12 = 0.
\]
On factorise immédiatement le trinôme par recherche de racines évidentes :
\[
(X - 2)(X - 6) = 0.
\]
Les solutions de cette équation sont $X_1 = 2$ et $X_2 = 6$.

\textbf{Conclusion :}\\
Les deux résistances valent :
\[
\boxed{R_1 = 2\,\Omega \quad \text{et} \quad R_2 = 6\,\Omega \quad (\text{à l'ordre près})}
\]
\textit{Vérification :} En série : $2 + 6 = 8\,\Omega$. En dérivation : $\frac{2 \times 6}{2 + 6} = \frac{12}{8} = 1{,}5\,\Omega$.
\end{correctionbox}

% ----------------------------------------------------
% EXERCICE 26
% ----------------------------------------------------
\subsection*{Exercice 26}
Un cycliste A a effectué un parcours de $55\text{ km}$ en $18\text{ min}$ de moins qu'un cycliste B, alors qu'il roule à une vitesse supérieure de $3\text{ km}\cdot\text{h}^{-1}$. Calculer la vitesse des deux cyclistes.

\begin{correctionbox}
\textbf{1. Choix des variables et mise en équation :}\\
Soit $v_B$ la vitesse du cycliste B, exprimée en $\text{km}\cdot\text{h}^{-1}$ ($v_B > 0$).\\
Le cycliste A roule à une vitesse supérieure de $3\text{ km}\cdot\text{h}^{-1}$, sa vitesse est donc :
\[
v_A = v_B + 3.
\]
La distance totale parcourue par chaque cycliste est $d = 55\text{ km}$.
Les durées des parcours (en heures) sont données par la relation fondamentale de la cinématique uniforme $t = \frac{d}{v}$ :
\[
t_A = \frac{55}{v_A} = \frac{55}{v_B + 3} \quad \text{et} \quad t_B = \frac{55}{v_B}.
\]
L'écart temporel entre les deux cyclistes est de $18\text{ min}$. Convertissons cette durée en heures :
\[
18\text{ min} = \frac{18}{60}\text{ h} = \frac{3}{10}\text{ h} = 0{,}3\text{ h}.
\]
Le cycliste A met moins de temps que le cycliste B, ce qui donne la relation :
\[
t_B - t_A = \frac{3}{10} \iff \frac{55}{v_B} - \frac{55}{v_B + 3} = \frac{3}{10}.
\]

\textbf{2. Réduction à une équation du second degré :}\\
Mettons le membre de gauche au même dénominateur :
\[
55 \left( \frac{(v_B + 3) - v_B}{v_B(v_B + 3)} \right) = \frac{3}{10} \iff \frac{55 \times 3}{v_B(v_B + 3)} = \frac{3}{10}.
\]
En simplifiant par $3$ (qui est non nul) de chaque côté de l'égalité :
\[
\frac{55}{v_B(v_B + 3)} = \frac{1}{10}.
\]
Le produit en croix donne :
\[
v_B(v_B + 3) = 55 \times 10 = 550 \iff v_B^2 + 3v_B - 550 = 0.
\]

\textbf{3. Résolution de l'équation du second degré :}\\
Calculons le discriminant du trinôme :
\[
\Delta = 3^2 - 4 \times 1 \times (-550) = 9 + 2200 = 2209.
\]
Calculons $\sqrt{2209}$ : puisque $40^2 = 1600$ et $50^2 = 2500$, le nombre cherché est compris entre $40$ et $50$. Comme son dernier chiffre est $9$, il ne peut s'agir que de $43$ ou $47$.
Or $47^2 = (50 - 3)^2 = 2500 - 300 + 9 = 2209$. Ainsi $\sqrt{\Delta} = 47$.

L'équation admet deux racines réelles :
\[
v_{B,1} = \frac{-3 - 47}{2} = -\frac{50}{2} = -25, \qquad v_{B,2} = \frac{-3 + 47}{2} = \frac{44}{2} = 22.
\]
Une vitesse étant une grandeur physique strictement positive, on écarte la racine négative $-25$. On a donc :
\[
v_B = 22\text{ km}\cdot\text{h}^{-1}.
\]
On en déduit la vitesse du cycliste A :
\[
v_A = v_B + 3 = 22 + 3 = 25\text{ km}\cdot\text{h}^{-1}.
\]

\textbf{Conclusion et vérification :}\\
\[
\boxed{v_A = 25\text{ km}\cdot\text{h}^{-1} \quad \text{et} \quad v_B = 22\text{ km}\cdot\text{h}^{-1}}
\]
\textit{Vérification :}
$t_B = \frac{55}{22}\text{ h} = \frac{5}{2}\text{ h} = 2\text{ h } 30\text{ min}$, et $t_A = \frac{55}{25}\text{ h} = \frac{11}{5}\text{ h} = 2{,}2\text{ h} = 2\text{ h } 12\text{ min}$.
L'écart est bien de $2\text{ h } 30\text{ min} - 2\text{ h } 12\text{ min} = 18\text{ min}$.
\end{correctionbox}

% ----------------------------------------------------
% EXERCICE 27
% ----------------------------------------------------
\subsection*{Exercice 27}
Soit $n \in \mathbb{N}$. Déterminer suivant les valeurs de l'entier $n$ le nombre de solutions dans $\mathbb{R}^2$ du système :
\[
\begin{cases}
x + y = n \\
xy = n
\end{cases}
\]

\begin{correctionbox}
\textbf{1. Propriété caractéristique des racines :}\\
Un couple de réels $(x, y) \in \mathbb{R}^2$ est solution du système symétrique si et seulement si $x$ et $y$ sont les racines réelles de l'équation polynomiale du second degré de variable $X$ :
\[
X^2 - SX + P = 0 \iff X^2 - nX + n = 0.
\]

\textbf{2. Étude du discriminant :}\\
Le discriminant de ce trinôme dépend du paramètre $n$ :
\[
\Delta(n) = (-n)^2 - 4 \times 1 \times n = n^2 - 4n = n(n - 4).
\]
Le nombre de solutions réelles de l'équation dépend directement du signe de $\Delta(n)$ sur $\mathbb{N}$ :

\begin{itemize}[leftmargin=*,label=\textbullet]
    \item \textbf{Pour $n = 0$ :}
    $\Delta(0) = 0$. Le trinôme devient $X^2 = 0$, qui admet l'unique racine double $X_0 = 0$.
    On a alors $x = y = 0$.
    Le système admet un \textbf{unique couple solution} : $(0\,;\,0)$.

    \item \textbf{Pour $n \in \{1, 2, 3\}$ :}
    Le tableau de signes de la fonction $t \mapsto t(t-4)$ montre que pour $t \in ]0, 4[$, $\Delta(t) < 0$ :
    \begin{itemize}
        \item Si $n = 1$, $\Delta(1) = 1(1 - 4) = -3 < 0$ ;
        \item Si $n = 2$, $\Delta(2) = 2(2 - 4) = -4 < 0$ ;
        \item Si $n = 3$, $\Delta(3) = 3(3 - 4) = -3 < 0$.
    \end{itemize}
    Puisque $\Delta(n) < 0$, l'équation $X^2 - nX + n = 0$ ne possède aucune solution réelle.
    Le système n'admet donc \textbf{aucune solution dans $\mathbb{R}^2$} ($0$ solution).

    \item \textbf{Pour $n = 4$ :}
    $\Delta(4) = 4(4 - 4) = 0$. L'équation $X^2 - 4X + 4 = 0 \iff (X - 2)^2 = 0$ admet une unique racine double $X_0 = 2$.
    On a alors $x = y = 2$.
    Le système admet un \textbf{unique couple solution} : $(2\,;\,2)$.

    \item \textbf{Pour $n \ge 5$ :}
    On a $n > 4$, donc $\Delta(n) = n(n - 4) > 0$.
    L'équation du second degré admet deux racines réelles distinctes :
    \[
    X_1 = \frac{n - \sqrt{n(n-4)}}{2} \quad \text{et} \quad X_2 = \frac{n + \sqrt{n(n-4)}}{2}.
    \]
    Comme $X_1 \ne X_2$, on obtient exactement deux couples ordonnés $(x, y)$ distincts en affectant ces deux racines à $(x, y)$ :
    \[
    (x, y) = (X_1, X_2) \quad \text{ou} \quad (x, y) = (X_2, X_1).
    \]
    Le système admet donc \textbf{exactement $2$ couples solutions distincts}.
\end{itemize}

\textbf{Synthèse des résultats :}
\begin{center}
\renewcommand{\arraystretch}{1.3}
\begin{tabular}{|c|c|l|}
\hline
\textbf{Valeur du paramètre $n \in \mathbb{N}$} & \textbf{Signe de $\Delta(n)$} & \textbf{Nombre de solutions du système} \\
\hline
$n \in \{1, 2, 3\}$ & $\Delta(n) < 0$ & \textbf{Aucune solution} ($0$ couple) \\
\hline
$n = 0$ ou $n = 4$ & $\Delta(n) = 0$ & \textbf{1 unique couple solution} : $(0,0)$ pour $n=0$, $(2,2)$ pour $n=4$ \\
\hline
$n \ge 5$ & $\Delta(n) > 0$ & \textbf{2 couples solutions distincts} \\
\hline
\end{tabular}
\end{center}
\end{correctionbox}

% ----------------------------------------------------
% EXERCICE 28
% ----------------------------------------------------
\subsection*{Exercice 28}
Déterminer les valeurs de $z$ telles que le système suivant admette des solutions dans $\mathbb{R}^2$ :
\[
\begin{cases}
x^2 + y^2 + z^2 = 4 \\
x + y + z = 3
\end{cases}
\]
Interpréter géométriquement ce résultat (on pourra s'aider de GeoGebra).

\begin{correctionbox}
\textbf{1. Approche analytique (trinôme du second degré en $x$) :}\\
Considérons le réel $z \in \mathbb{R}$ comme un paramètre fixé.
Le système équivaut à trouver l'existence d'au moins un couple de réels $(x, y)$ vérifiant :
\[
\begin{cases}
x^2 + y^2 = 4 - z^2 \\
x + y = 3 - z
\end{cases}
\]
De la seconde équation, exprimons $y$ en fonction de $x$ et $z$ :
\[
y = (3 - z) - x.
\]
Injectons cette expression dans la première équation :
\[
x^2 + \bigl((3 - z) - x\bigr)^2 = 4 - z^2.
\]
Développons le terme quadratique :
\[
x^2 + \bigl((3 - z)^2 - 2(3 - z)x + x^2\bigr) = 4 - z^2
\]
\[
\iff 2x^2 - 2(3 - z)x + (9 - 6z + z^2) - (4 - z^2) = 0
\]
\[
\iff 2x^2 + 2(z - 3)x + (2z^2 - 6z + 5) = 0.
\]
Cette équation d'inconnue $x$ est une équation polynomiale du second degré (le coefficient dominant est $2 \neq 0$).
Elle admet des solutions réelles si et seulement si son discriminant $\Delta_x$ est supérieur ou égal à zéro.

Calculons $\Delta_x$ (ou le discriminant réduit $\Delta'_x = \frac{\Delta_x}{4}$) :
\[
\begin{aligned}
\Delta'_x = \left(\frac{b}{2}\right)^2 - ac &= (z - 3)^2 - 2(2z^2 - 6z + 5) \\
&= (z^2 - 6z + 9) - (4z^2 - 12z + 10) \\
&= -3z^2 + 6z - 1.
\end{aligned}
\]
Le système admet au moins une solution si et seulement si :
\[
-3z^2 + 6z - 1 \ge 0 \iff 3z^2 - 6z + 1 \le 0.
\]

\textbf{2. Résolution de l'inéquation en $z$ :}\\
Étudions le trinôme $T(z) = 3z^2 - 6z + 1$ :
\[
\Delta_z = (-6)^2 - 4 \times 3 \times 1 = 36 - 12 = 24 = (2\sqrt{6})^2 > 0.
\]
Le trinôme admet deux racines réelles distinctes :
\[
z_1 = \frac{6 - 2\sqrt{6}}{2 \times 3} = 1 - \frac{\sqrt{6}}{3}, \qquad z_2 = \frac{6 + 2\sqrt{6}}{2 \times 3} = 1 + \frac{\sqrt{6}}{3}.
\]
Le coefficient de $z^2$ étant strictement positif ($a = 3 > 0$), le trinôme $3z^2 - 6z + 1$ est négatif ou nul entre ses racines.

\textbf{Conclusion pour $z$ :}\\
L'ensemble des valeurs de $z$ pour lesquelles le système admet des solutions est l'intervalle fermé :
\[
\boxed{z \in \left[ 1 - \frac{\sqrt{6}}{3} \,;\, 1 + \frac{\sqrt{6}}{3} \right]}
\]
\textit{Remarque :} On a $1 - \frac{\sqrt{6}}{3} \approx 1 - 0{,}816 = 0{,}184$ et $1 + \frac{\sqrt{6}}{3} \approx 1{,}816$. Comme $|z| < 2$, on a bien $4 - z^2 > 0$.

\textbf{3. Méthode alternative : Inégalité de Cauchy-Schwarz / Convexité}\\
Pour tous réels $x$ et $y$, on a $(x - y)^2 \ge 0 \iff 2(x^2 + y^2) \ge (x + y)^2$.\\
En substituant $x^2 + y^2 = 4 - z^2$ et $x + y = 3 - z$, on obtient immédiatement :
\[
2(4 - z^2) \ge (3 - z)^2 \iff 8 - 2z^2 \ge 9 - 6z + z^2 \iff 3z^2 - 6z + 1 \le 0,
\]
ce qui confirme instantanément l'encadrement ci-dessus.

\textbf{4. Interprétation géométrique dans l'espace $\mathbb{R}^3$ :}\\
Muni d'un repère orthonormé $(O; \vec{i}, \vec{j}, \vec{k})$ :
\begin{itemize}[leftmargin=*,label=\textbullet]
    \item L'équation $x^2 + y^2 + z^2 = 4 = 2^2$ est celle de la \textbf{sphère} $\mathcal{S}$ de centre l'origine $O(0,0,0)$ et de rayon $R = 2$.
    \item L'équation $x + y + z = 3$ est celle d'un \textbf{plan} $\mathcal{P}$ de vecteur normal $\vec{n}(1, 1, 1)$.
    \item La distance du centre $O$ au plan $\mathcal{P}$ vaut :
    \[
    d(O, \mathcal{P}) = \frac{|0 + 0 + 0 - 3|}{\sqrt{1^2 + 1^2 + 1^2}} = \frac{3}{\sqrt{3}} = \sqrt{3}.
    \]
    Comme $d(O, \mathcal{P}) = \sqrt{3} \approx 1{,}732 < R = 2$, le plan coupe la sphère selon un \textbf{cercle de l'espace} $\mathcal{C}$, de rayon :
    \[
    r = \sqrt{R^2 - d^2} = \sqrt{4 - 3} = 1.
    \]
    Le centre de ce cercle est le projeté orthogonal $\Omega$ de $O$ sur $\mathcal{P}$, c'est-à-dire le point $\Omega(1, 1, 1)$.
    Les valeurs de $z$ correspondent à la projection orthogonale de ce cercle $\mathcal{C}$ sur l'axe vertical $(Oz)$. L'altitude minimale et l'altitude maximale des points du cercle sont précisément $z_{\min} = 1 - \frac{\sqrt{6}}{3}$ et $z_{\max} = 1 + \frac{\sqrt{6}}{3}$.
\end{itemize}
\end{correctionbox}

% ----------------------------------------------------
% EXERCICE 29
% ----------------------------------------------------
\subsection*{Exercice 29}
Résoudre dans $\mathbb{R}^2$ le système :
\[
\begin{cases}
x^2 + xy + y^2 = 7 \\
(x^2 + 1)(y^2 + 1) + 2xy = 14
\end{cases}
\]
Vérifier vos solutions à l'aide de GeoGebra.

\begin{correctionbox}
\textbf{1. Transformation et combinaison algébrique :}\\
Développons la seconde équation du système :
\[
(x^2 + 1)(y^2 + 1) + 2xy = 14 \iff x^2 y^2 + x^2 + y^2 + 1 + 2xy = 14.
\]
Remarquons que :
\[
x^2 + y^2 + 2xy = (x^2 + xy + y^2) + xy.
\]
En réinjectant cette écriture dans l'équation développée :
\[
(xy)^2 + (x^2 + xy + y^2) + xy + 1 = 14.
\]
Or, la première équation nous donne $x^2 + xy + y^2 = 7$. Remplaçons cette somme par $7$ :
\[
(xy)^2 + 7 + xy + 1 = 14 \iff (xy)^2 + xy + 8 = 14 \iff (xy)^2 + xy - 6 = 0.
\]

\textbf{2. Détermination des valeurs du produit $Z = xy$ :}\\
Posons la variable auxiliaire $Z = xy$. L'équation devient :
\[
Z^2 + Z - 6 = 0.
\]
Discriminant : $\Delta = 1^2 - 4 \times 1 \times (-6) = 25 = 5^2$.
Les solutions sont :
\[
Z_1 = \frac{-1 - 5}{2} = -3 \quad \text{et} \quad Z_2 = \frac{-1 + 5}{2} = 2.
\]
Nous distinguons donc deux cas disjoints.

\textbf{Premier cas : $xy = 2$}\\
La première équation donne $x^2 + 2 + y^2 = 7 \iff x^2 + y^2 = 5$.\\
Calculons la somme $S = x + y$ :
\[
(x + y)^2 = x^2 + y^2 + 2xy = 5 + 2 \times 2 = 9 \implies x + y = 3 \quad \text{ou} \quad x + y = -3.
\]
\begin{itemize}[leftmargin=*,label=\textbullet]
    \item Si $x + y = 3$ et $xy = 2$ : $x$ et $y$ sont les racines de $X^2 - 3X + 2 = 0 \iff (X - 1)(X - 2) = 0$.
    On obtient les couples : $(1\,;\,2)$ et $(2\,;\,1)$.
    \item Si $x + y = -3$ et $xy = 2$ : $x$ et $y$ sont les racines de $X^2 + 3X + 2 = 0 \iff (X + 1)(X + 2) = 0$.
    On obtient les couples : $(-1\,;\,-2)$ et $(-2\,;\,-1)$.
\end{itemize}

\textbf{Deuxième cas : $xy = -3$}\\
La première équation donne $x^2 - 3 + y^2 = 7 \iff x^2 + y^2 = 10$.\\
Puisque $xy = -3 \ne 0$, on a $y = -\frac{3}{x}$. Substituons dans $x^2 + y^2 = 10$ :
\[
x^2 + \frac{9}{x^2} = 10 \iff x^4 - 10x^2 + 9 = 0.
\]
En posant $u = x^2 \ge 0$, on a l'équation $u^2 - 10u + 9 = 0 \iff (u - 1)(u - 9) = 0$.\\
D'où $x^2 = 1$ ou $x^2 = 9$.
\begin{itemize}[leftmargin=*,label=\textbullet]
    \item Si $x^2 = 1 \iff x = \pm 1$ :
    \begin{itemize}
        \item Pour $x = 1$ : $y = -\frac{3}{1} = -3 \implies (1\,;\,-3)$ ;
        \item Pour $x = -1$ : $y = -\frac{3}{-1} = 3 \implies (-1\,;\,3)$.
    \end{itemize}
    \item Si $x^2 = 9 \iff x = \pm 3$ :
    \begin{itemize}
        \item Pour $x = 3$ : $y = -\frac{3}{3} = -1 \implies (3\,;\,-1)$ ;
        \item Pour $x = -3$ : $y = -\frac{3}{-3} = 1 \implies (-3\,;\,1)$.
    \end{itemize}
\end{itemize}

\textbf{Conclusion :}\\
Le système admet exactement \textbf{8 couples de solutions} dans $\mathbb{R}^2$ :
\[
\boxed{\mathcal{S} = \bigl\{ (1\,;\,2),\, (2\,;\,1),\, (-1\,;\,-2),\, (-2\,;\,-1),\, (1\,;\,-3),\, (-1\,;\,3),\, (3\,;\,-1),\, (-3\,;\,1) \bigr\}}
\]
\textit{Vérification GeoGebra :} En traçant les courbes implicites d'équations $x^2 + xy + y^2 = 7$ (qui est une ellipse inclinée) et $(x^2+1)(y^2+1)+2xy = 14$ (courbe algébrique de degré 4), on observe exactement $8$ points d'intersection qui coïncident avec les coordonnées déterminées.
\end{correctionbox}


% ==============================================================================
% SECTION 2 : AUTOUR DU DEGRÉ 3
% ==============================================================================
\section{Autour du degré 3}

% ----------------------------------------------------
% EXERCICE 30
% ----------------------------------------------------
\subsection*{Exercice 30 (Identification polynomiale de degré 2)}
Soient $a, b, c, a', b'$ et $c'$ des réels et $f$ et $g$ les fonctions définies sur $\mathbb{R}$ par :
\[
f(x) = ax^2 + bx + c \quad \text{et} \quad g(x) = a'x^2 + b'x + c'.
\]
\begin{enumerate}
    \item Montrer que : La fonction $f$ est la fonction nulle sur $\mathbb{R} \iff a = b = c = 0$.
    \item Montrer que : $f = g \iff \begin{cases} a = a' \\ b = b' \\ c = c' \end{cases}$.
\end{enumerate}

\begin{correctionbox}
\begin{enumerate}
    \item \textbf{Étude de l'équivalence pour la fonction nulle :}\\
    $\bullet$ \textbf{Sens réciproque ($\Longleftarrow$) :} Si $a = b = c = 0$, alors pour tout $x \in \mathbb{R}$, on a :
    \[
    f(x) = 0 \times x^2 + 0 \times x + 0 = 0.
    \]
    La fonction $f$ est donc identiquement nulle sur $\mathbb{R}$.

    $\bullet$ \textbf{Sens direct ($\Longrightarrow$) :} Supposons que $f$ est la fonction nulle sur $\mathbb{R}$, c'est-à-dire :
    \[
    \forall x \in \mathbb{R}, \quad ax^2 + bx + c = 0.
    \]
    Évaluons cette égalité en trois points particuliers judicieusement choisis :
    \begin{itemize}[leftmargin=*,label=\textbullet]
        \item Pour $x = 0$ :
        \[
        f(0) = a(0)^2 + b(0) + c = 0 \implies c = 0.
        \]
        L'expression se simplifie donc en $f(x) = ax^2 + bx$.
        \item Pour $x = 1$ :
        \[
        f(1) = a(1)^2 + b(1) = 0 \implies a + b = 0 \quad (L_1).
        \]
        \item Pour $x = -1$ :
        \[
        f(-1) = a(-1)^2 + b(-1) = 0 \implies a - b = 0 \quad (L_2).
        \]
    \end{itemize}
    En effectuant la demi-somme et la demi-différence des relations $(L_1)$ et $(L_2)$ :
    \[
    (L_1) + (L_2) \implies 2a = 0 \implies a = 0,
    \]
    \[
    (L_1) - (L_2) \implies 2b = 0 \implies b = 0.
    \]
    On a donc bien établi que $a = 0$, $b = 0$ et $c = 0$.

    \textit{Autre démonstration par l'absurde :} Si $(a, b) \neq (0, 0)$, alors $f$ est une fonction polynomiale de degré $1$ ou $2$, elle admet donc au plus $2$ racines réelles distinctes. Elle ne peut donc pas s'annuler sur tout $\mathbb{R}$. Par conséquent, $a = b = c = 0$.

    \item \textbf{Égalité de deux fonctions polynomiales :}\\
    Deux fonctions $f$ et $g$ sont égales sur $\mathbb{R}$ si et seulement si leur différence $f - g$ est la fonction nulle sur $\mathbb{R}$ :
    \[
    f = g \iff \forall x \in \mathbb{R}, \quad (f - g)(x) = 0.
    \]
    Or, pour tout $x \in \mathbb{R}$ :
    \[
    (f - g)(x) = (ax^2 + bx + c) - (a'x^2 + b'x + c') = (a - a')x^2 + (b - b')x + (c - c').
    \]
    La fonction $f - g$ est une fonction polynôme du second degré de coefficients $(a - a')$, $(b - b')$ et $(c - c')$. D'après le résultat démontré à la question 1, elle est nulle sur $\mathbb{R}$ si et seulement si tous ses coefficients sont nuls :
    \[
    \begin{cases}
    a - a' = 0 \\
    b - b' = 0 \\
    c - c' = 0
    \end{cases}
    \iff
    \begin{cases}
    a = a' \\
    b = b' \\
    c = c'
    \end{cases}
    \]
    Ce résultat fondamental justifie rigoureusement la \textbf{méthode d'identification des coefficients} pour les polynômes.
\end{enumerate}
\end{correctionbox}

% ----------------------------------------------------
% EXERCICE 31
% ----------------------------------------------------
\subsection*{Exercice 31 (Factorisation d'un polynôme de degré 3)}
Soient $a, b, c$ et $d$ des réels ($a \neq 0$).
Soit $f$ la fonction polynomiale définie sur $\mathbb{R}$ par $f(x) = ax^3 + bx^2 + cx + d$. Soit $r$ une racine de $f$.
\begin{enumerate}
    \item Montrer qu'il existe une fonction polynomiale $g$ de degré 2 telle que pour tout $x \in \mathbb{R}$, $f(x) = (x - r)g(x)$.
    \item Application : Résoudre dans $\mathbb{R}$ l'équation $3x^3 - 2x^2 + 5x - 6 = 0$.
\end{enumerate}

\begin{correctionbox}
\begin{enumerate}
    \item \textbf{Démonstration de l'existence de la factorisation :}\\
    Puisque $r$ est une racine de $f$, on a $f(r) = 0$. On peut donc écrire, pour tout réel $x$ :
    \[
    f(x) = f(x) - f(r) = (ax^3 + bx^2 + cx + d) - (ar^3 + br^2 + cr + d).
    \]
    En regroupant les termes selon les puissances de même degré :
    \[
    f(x) = a(x^3 - r^3) + b(x^2 - r^2) + c(x - r).
    \]
    Utilisons les identités remarquables de factorisation :
    \[
    x^2 - r^2 = (x - r)(x + r) \quad \text{et} \quad x^3 - r^3 = (x - r)(x^2 + xr + r^2).
    \]
    Le terme $(x - r)$ apparaît en facteur commun :
    \[
    \begin{aligned}
    f(x) &= a(x - r)(x^2 + xr + r^2) + b(x - r)(x + r) + c(x - r) \\
    &= (x - r) \Bigl[ a(x^2 + xr + r^2) + b(x + r) + c \Bigr] \\
    &= (x - r) \Bigl[ ax^2 + (ar + b)x + (ar^2 + br + c) \Bigr].
    \end{aligned}
    \]
    Posons alors la fonction $g$ définie sur $\mathbb{R}$ par :
    \[
    g(x) = ax^2 + (ar + b)x + (ar^2 + br + c).
    \]
    Comme $a \ne 0$, $g$ est bien une fonction polynomiale de degré 2, et on a :
    \[
    \forall x \in \mathbb{R}, \quad f(x) = (x - r)g(x).
    \]

    \item \textbf{Application : Résolution de $3x^3 - 2x^2 + 5x - 6 = 0$ :}\\
    Posons $P(x) = 3x^3 - 2x^2 + 5x - 6$.
    \begin{itemize}[leftmargin=*,label=\textbullet]
        \item \textbf{Recherche d'une racine évidente :}
        Testons $x = 1$ :
        \[
        P(1) = 3(1)^3 - 2(1)^2 + 5(1) - 6 = 3 - 2 + 5 - 6 = 0.
        \]
        Donc $r = 1$ est une racine évidente de $P$.

        \item \textbf{Factorisation par $(x - 1)$ :}
        D'après la question 1, il existe des réels $A, B, C$ tels que :
        \[
        P(x) = (x - 1)(Ax^2 + Bx + C), \quad \text{avec } A = 3.
        \]
        Développons cette expression :
        \[
        (x - 1)(3x^2 + Bx + C) = 3x^3 + (B - 3)x^2 + (C - B)x - C.
        \]
        Par identification avec $3x^3 - 2x^2 + 5x - 6$ (valable d'après l'unicité prouvée à l'exercice 30) :
        \[
        \begin{cases}
        B - 3 = -2 \implies B = 1 \\
        C - B = 5 \implies C - 1 = 5 \implies C = 6 \\
        -C = -6 \implies C = 6
        \end{cases}
        \]
        Le système est parfaitement cohérent, d'où :
        \[
        P(x) = (x - 1)(3x^2 + x + 6).
        \]

        \item \textbf{Résolution de l'équation $P(x) = 0$ :}
        Un produit de facteurs est nul si et seulement si l'un au moins des facteurs est nul :
        \[
        (x - 1)(3x^2 + x + 6) = 0 \iff x - 1 = 0 \quad \text{ou} \quad 3x^2 + x + 6 = 0.
        \]
        Étudions le trinôme du second degré $3x^2 + x + 6$ :
        \[
        \Delta = 1^2 - 4 \times 3 \times 6 = 1 - 72 = -71 < 0.
        \]
        Puisque $\Delta < 0$, ce trinôme n'admet aucune racine réelle.
    \end{itemize}

    \textbf{Conclusion :}\\
    L'équation admet une unique solution dans $\mathbb{R}$ :
    \[
    \boxed{\mathcal{S} = \{ 1 \}}
    \]
\end{enumerate}
\end{correctionbox}

% ----------------------------------------------------
% EXERCICE 32
% ----------------------------------------------------
\subsection*{Exercice 32}
Montrer que le nombre réel suivant est un entier naturel :
\[
A = \sqrt[3]{3\sqrt{21} + 8} - \sqrt[3]{3\sqrt{21} - 8} \in \mathbb{N}.
\]
\textit{Indication :} On pourra montrer que $A$ est une solution d'une équation de degré 3.

\begin{correctionbox}
\textbf{1. Analyse des radicaux et pose des variables auxiliaires :}\\
Posons :
\[
u = \sqrt[3]{3\sqrt{21} + 8} \quad \text{et} \quad v = \sqrt[3]{3\sqrt{21} - 8}.
\]
Remarquons que $3\sqrt{21} = \sqrt{9 \times 21} = \sqrt{189} > \sqrt{64} = 8$, donc $3\sqrt{21} - 8 > 0$. Les réels $u$ et $v$ sont donc strictement positifs, et comme $3\sqrt{21} + 8 > 3\sqrt{21} - 8$, on a $u > v$, de sorte que $A = u - v > 0$.

\textbf{2. Calcul du produit $uv$ :}\\
En regroupant sous un même radical cubique :
\[
uv = \sqrt[3]{(3\sqrt{21} + 8)(3\sqrt{21} - 8)}.
\]
Grâce à l'identité remarquable $(a+b)(a-b) = a^2 - b^2$ :
\[
(3\sqrt{21} + 8)(3\sqrt{21} - 8) = (3\sqrt{21})^2 - 8^2 = 9 \times 21 - 64 = 189 - 64 = 125.
\]
Comme $125 = 5^3$, on en déduit que :
\[
uv = \sqrt[3]{125} = 5.
\]

\textbf{3. Établissement de l'équation satisfaite par $A$ :}\\
Calculons $A^3 = (u - v)^3$ en utilisant le développement du binôme de Newton sous la forme factorisée :
\[
(u - v)^3 = u^3 - 3u^2 v + 3uv^2 - v^3 = (u^3 - v^3) - 3uv(u - v).
\]
Calculons chaque bloc séparément :
\begin{itemize}[leftmargin=*,label=\textbullet]
    \item $u^3 - v^3 = (3\sqrt{21} + 8) - (3\sqrt{21} - 8) = 8 + 8 = 16$ ;
    \item $uv = 5$ ;
    \item $u - v = A$.
\end{itemize}
En injectant ces données dans l'égalité cubique :
\[
A^3 = 16 - 3 \times 5 \times A = 16 - 15A.
\]
Par conséquent, le réel $A$ est une solution de l'équation polynomiale du troisième degré :
\[
A^3 + 15A - 16 = 0.
\]

\textbf{4. Résolution de l'équation $X^3 + 15X - 16 = 0$ dans $\mathbb{R}$ :}\\
Soit $P(X) = X^3 + 15X - 16$.
\begin{itemize}[leftmargin=*,label=\textbullet]
    \item $X = 1$ est une racine évidente de $P$ car $1^3 + 15(1) - 16 = 16 - 16 = 0$.
    \item Factorisons $P(X)$ par $(X - 1)$ :
    \[
    X^3 + 15X - 16 = (X - 1)(X^2 + X + 16).
    \]
    \textit{Vérification du développement :}
    $(X - 1)(X^2 + X + 16) = X^3 + X^2 + 16X - X^2 - X - 16 = X^3 + 15X - 16$.
    \item Étudions le facteur du second degré $Q(X) = X^2 + X + 16$ :
    \[
    \Delta_Q = 1^2 - 4 \times 1 \times 16 = 1 - 64 = -63 < 0.
    \]
    Le discriminant étant strictement négatif, ce trinôme ne s'annule pour aucune valeur réelle et reste strictement positif sur $\mathbb{R}$.
\end{itemize}
L'équation $X^3 + 15X - 16 = 0$ admet donc \textbf{une unique racine réelle}, qui est $X = 1$.

\textbf{Conclusion :}\\
Puisque $A$ est un réel vérifiant cette équation, on a nécessairement :
\[
\boxed{A = 1 \in \mathbb{N}}
\]
\end{correctionbox}

% ----------------------------------------------------
% EXERCICE 33
% ----------------------------------------------------
\subsection*{Exercice 33}
\begin{enumerate}
    \item Déterminer un polynôme $P$ de degré 3 vérifiant les deux conditions suivantes :
    \begin{enumerate}[label=(\alph*)]
        \item $\forall x \in \mathbb{R}, \quad P(x+1) - P(x) = x(x+1)$
        \item $P(0) = 0$
    \end{enumerate}
    \item Pour tout entier naturel $n \ge 1$, on pose $S(n) = \displaystyle\sum_{k=1}^n k(k+1)$.\\
    Montrer que $S(n) = \dfrac{n(n+1)(n+2)}{3}$.
\end{enumerate}

\begin{correctionbox}
\begin{enumerate}
    \item \textbf{Recherche du polynôme $P$ :}\\
    Puisque $P$ est de degré $3$, il s'écrit sous la forme générale :
    \[
    P(x) = ax^3 + bx^2 + cx + d, \quad \text{avec } a \neq 0.
    \]
    \begin{itemize}[leftmargin=*,label=\textbullet]
        \item La condition (b) donne directement :
        \[
        P(0) = d = 0.
        \]
        Ainsi, $P(x) = ax^3 + bx^2 + cx$.
        \item Développons la différence finie $P(x+1) - P(x)$ :
        \[
        P(x+1) - P(x) = a\bigl((x+1)^3 - x^3\bigr) + b\bigl((x+1)^2 - x^2\bigr) + c\bigl((x+1) - x\bigr).
        \]
        Calculons chaque terme :
        \[
        (x+1)^3 - x^3 = (x^3 + 3x^2 + 3x + 1) - x^3 = 3x^2 + 3x + 1,
        \]
        \[
        (x+1)^2 - x^2 = (x^2 + 2x + 1) - x^2 = 2x + 1,
        \]
        \[
        (x+1) - x = 1.
        \]
        En regroupant par puissances de $x$ :
        \[
        \begin{aligned}
        P(x+1) - P(x) &= a(3x^2 + 3x + 1) + b(2x + 1) + c \\
        &= 3ax^2 + (3a + 2b)x + (a + b + c).
        \end{aligned}
        \]
        \item Or, la condition (a) impose que pour tout réel $x$ :
        \[
        P(x+1) - P(x) = x(x+1) = x^2 + x.
        \]
        \item Par identification des coefficients (valable sur les polynômes réels) :
        \[
        \begin{cases}
        3a = 1 \\[1mm]
        3a + 2b = 1 \\[1mm]
        a + b + c = 0
        \end{cases}
        \iff
        \begin{cases}
        a = \dfrac{1}{3} \\[2mm]
        1 + 2b = 1 \implies b = 0 \\[1mm]
        \dfrac{1}{3} + 0 + c = 0 \implies c = -\dfrac{1}{3}
        \end{cases}
        \]
    \end{itemize}
    Le polynôme cherché est donc :
    \[
    \boxed{P(x) = \frac{1}{3}x^3 - \frac{1}{3}x = \frac{1}{3}x(x^2 - 1) = \frac{1}{3}(x - 1)x(x + 1)}
    \]

    \item \textbf{Calcul de la somme $S(n)$ par télescopage :}\\
    Pour tout entier $k \ge 1$, la condition vérifiée par le polynôme $P$ donne :
    \[
    k(k+1) = P(k+1) - P(k).
    \]
    En sommant cette identité pour $k$ allant de $1$ à $n$ :
    \[
    S(n) = \sum_{k=1}^n k(k+1) = \sum_{k=1}^n \bigl( P(k+1) - P(k) \bigr).
    \]
    La somme est télescopique : tous les termes intermédiaires s'annulent deux à deux :
    \[
    S(n) = \bigl( P(2) - P(1) \bigr) + \bigl( P(3) - P(2) \bigr) + \dots + \bigl( P(n+1) - P(n) \bigr) = P(n+1) - P(1).
    \]
    Or $P(1) = \frac{1}{3}(1)(1-1)(1+1) = 0$.\\
    Il vient donc simplement :
    \[
    S(n) = P(n+1).
    \]
    Évaluons $P(n+1)$ à l'aide de la forme factorisée $P(x) = \frac{1}{3}(x-1)x(x+1)$ :
    \[
    S(n) = \frac{1}{3}\bigl((n+1) - 1\bigr)(n+1)\bigl((n+1) + 1\bigr) = \frac{1}{3}n(n+1)(n+2).
    \]
    On a donc bien démontré que pour tout $n \in \mathbb{N}^*$ :
    \[
    \boxed{S(n) = \frac{n(n+1)(n+2)}{3}}
    \]
    \textit{Remarque culturelle :} On retrouve bien la somme classique $\sum_{k=1}^n (k^2+k) = \frac{n(n+1)(2n+1)}{6} + \frac{n(n+1)}{2} = \frac{n(n+1)}{6}[2n+1+3] = \frac{n(n+1)(2n+4)}{6} = \frac{n(n+1)(n+2)}{3}$.
\end{enumerate}
\end{correctionbox}

% ----------------------------------------------------
% EXERCICE 34
% ----------------------------------------------------
\subsection*{Exercice 34}
Soit $n$ un entier naturel non nul. On considère le polynôme $P$ défini sur $\mathbb{R}$ par :
\[
P(x) = x^3 - 3nx^2 + (3n^2 - 1)x - n(n^2 - 1).
\]
Déterminer les racines de $P$.

\begin{correctionbox}
\textbf{1. Recherche d'une racine remarquable en fonction de $n$ :}\\
Observons la structure des coefficients : le terme constant est $-n(n^2 - 1) = -n^3 + n$.
Évaluons le polynôme en $x = n$ :
\[
\begin{aligned}
P(n) &= n^3 - 3n(n^2) + (3n^2 - 1)n - n(n^2 - 1) \\
&= n^3 - 3n^3 + 3n^3 - n - n^3 + n \\
&= (1 - 3 + 3 - 1)n^3 + (-1 + 1)n \\
&= 0.
\end{aligned}
\]
Ainsi, pour tout entier $n \ge 1$, $x = n$ est une racine de $P$.

\textbf{2. Factorisation de $P(x)$ par $(x - n)$ :}\\
D'après le cours sur la division euclidienne et la factorisation polynomiale, il existe un trinôme $Q(x) = x^2 + Bx + C$ tel que :
\[
P(x) = (x - n)(x^2 + Bx + C).
\]
Développons ce produit :
\[
(x - n)(x^2 + Bx + C) = x^3 + (B - n)x^2 + (C - nB)x - nC.
\]
En identifiant les coefficients avec $P(x) = x^3 - 3nx^2 + (3n^2 - 1)x - n(n^2 - 1)$ :
\begin{itemize}[leftmargin=*,label=\textbullet]
    \item Coefficient de $x^2$ : $B - n = -3n \implies B = -2n$ ;
    \item Terme constant : $-nC = -n(n^2 - 1) \implies C = n^2 - 1$ ;
    \item Vérification sur le terme en $x$ : $C - nB = (n^2 - 1) - n(-2n) = n^2 - 1 + 2n^2 = 3n^2 - 1$, ce qui concorde rigoureusement.
\end{itemize}
Le polynôme se factorise donc sous la forme :
\[
P(x) = (x - n)\bigl(x^2 - 2nx + (n^2 - 1)\bigr).
\]

\textbf{3. Résolution du trinôme du second degré :}\\
Considérons le trinôme $Q(x) = x^2 - 2nx + (n^2 - 1)$.
On reconnaît immédiatement le début d'une identité remarquable (forme canonique) :
\[
Q(x) = (x^2 - 2nx + n^2) - 1 = (x - n)^2 - 1^2.
\]
D'après l'identité remarquable $A^2 - B^2 = (A - B)(A + B)$ :
\[
Q(x) = \bigl((x - n) - 1\bigr)\bigl((x - n) + 1\bigr) = (x - n - 1)(x - n + 1).
\]
La factorisation complète de $P(x)$ est donc :
\[
P(x) = (x - (n - 1))(x - n)(x - (n + 1)).
\]

\textbf{Conclusion :}\\
Le polynôme $P$ admet trois racines réelles distinctes, qui sont trois entiers consécutifs :
\[
\boxed{x_1 = n - 1, \qquad x_2 = n, \qquad x_3 = n + 1}
\]
\end{correctionbox}


% ==============================================================================
% SECTION 3 : DIVERS
% ==============================================================================
\section{Exercices divers et approfondissements}

% ----------------------------------------------------
% EXERCICE 35
% ----------------------------------------------------
\subsection*{Exercice 35}
Soit $ABC$ un triangle tel que :
\[
AB^2 + AC^2 + BC^2 = AB \cdot AC + AC \cdot BC + AB \cdot BC.
\]
Justifier que $ABC$ est un triangle équilatéral.

\begin{correctionbox}
\textbf{1. Choix des notations :}\\
Pour simplifier les écritures, posons les longueurs des trois côtés du triangle :
\[
a = BC, \quad b = AC, \quad c = AB.
\]
Puisque $a, b, c$ sont des longueurs de segments d'un vrai triangle, ce sont des réels strictement positifs.
L'égalité de l'énoncé s'écrit alors :
\[
c^2 + b^2 + a^2 = c \cdot b + b \cdot a + c \cdot a \iff a^2 + b^2 + c^2 - ab - bc - ca = 0.
\]

\textbf{2. Transformation en somme de carrés :}\\
Multiplions les deux membres de cette égalité par $2$ :
\[
2a^2 + 2b^2 + 2c^2 - 2ab - 2bc - 2ca = 0.
\]
Décomposons chaque terme $2x^2$ sous la forme $x^2 + x^2$ et regroupons astucieusement les termes pour faire apparaître trois identités remarquables :
\[
(a^2 - 2ab + b^2) + (b^2 - 2bc + c^2) + (c^2 - 2ca + a^2) = 0,
\]
ce qui s'écrit sous forme condensée :
\[
(a - b)^2 + (b - c)^2 + (c - a)^2 = 0.
\]

\textbf{3. Conclusion sur les longueurs :}\\
Dans l'ensemble des réels $\mathbb{R}$, le carré de tout nombre est positif ou nul :
\[
\forall (a, b, c) \in \mathbb{R}^3, \quad (a - b)^2 \ge 0, \quad (b - c)^2 \ge 0, \quad (c - a)^2 \ge 0.
\]
Une somme de trois réels positifs ou nuls est nulle si et seulement si chacun des termes est nul simultanément :
\[
\begin{cases}
(a - b)^2 = 0 \\
(b - c)^2 = 0 \\
(c - a)^2 = 0
\end{cases}
\iff
\begin{cases}
a - b = 0 \\
b - c = 0 \\
c - a = 0
\end{cases}
\iff
\begin{cases}
a = b \\
b = c \\
c = a
\end{cases}
\iff a = b = c.
\]
Ainsi, on a $AB = BC = CA$. Le triangle $ABC$ a ses trois côtés de même longueur :
\[
\boxed{\text{Le triangle } ABC \text{ est équilatéral.}}
\]
\end{correctionbox}

% ----------------------------------------------------
% EXERCICE 36
% ----------------------------------------------------
\subsection*{Exercice 36}
Soit $n$ un entier naturel non nul. Montrer qu'il existe un unique entier naturel $k$ tel que :
\[
\frac{k(k+1)}{2} \le n < \frac{(k+1)(k+2)}{2}.
\]

\begin{correctionbox}
\textbf{Méthode 1 : Résolution analytique par encadrement du second degré}\\
Considérons la fonction $f$ définie sur $[0, +\infty[$ par :
\[
f(x) = \frac{x(x+1)}{2} = \frac{1}{2}x^2 + \frac{1}{2}x.
\]
$f$ est une fonction polynôme du second degré de dérivée $f'(x) = x + \frac{1}{2} > 0$ sur $[0, +\infty[$.
Elle est donc continue et strictement croissante sur $[0, +\infty[$, avec $f(0) = 0$ et $\lim_{x \to +\infty} f(x) = +\infty$.

D'après le théorème de la bijection (corollaire du TVI), pour tout $n \in \mathbb{N}^*$, il existe un unique réel $x_0 \in ]0, +\infty[$ tel que $f(x_0) = n$ :
\[
\frac{x_0(x_0 + 1)}{2} = n \iff x_0^2 + x_0 - 2n = 0.
\]
Le discriminant de cette équation du second degré est :
\[
\Delta = 1^2 - 4(1)(-2n) = 1 + 8n > 0.
\]
L'équation admet deux racines réelles de signes contraires. L'unique racine strictement positive est :
\[
x_0 = \frac{-1 + \sqrt{1 + 8n}}{2}.
\]
Puisque la fonction $f$ est strictement croissante sur $[0, +\infty[$, l'encadrement $f(k) \le n < f(k+1)$ équivaut à :
\[
f(k) \le f(x_0) < f(k+1) \iff k \le x_0 < k+1.
\]
Cette double inégalité caractérise de manière unique la \textbf{partie entière} du réel $x_0$ :
\[
k = \lfloor x_0 \rfloor = \left\lfloor \frac{\sqrt{8n + 1} - 1}{2} \right\rfloor.
\]
Comme $n \ge 1$, on a $8n + 1 \ge 9$, d'où $\sqrt{8n + 1} \ge 3$, ce qui garantit $x_0 \ge \frac{3 - 1}{2} = 1$.
Ainsi, $k = \lfloor x_0 \rfloor$ est un entier naturel bien défini, supérieur ou égal à $1$, et son existence ainsi que son unicité sont rigoureusement établies.

\textbf{Méthode 2 : Partition de $\mathbb{N}^*$ par les nombres triangulaires}\\
Notons $T_m = \frac{m(m+1)}{2}$ le $m$-ième nombre triangulaire.
La suite $(T_m)_{m \in \mathbb{N}}$ est strictement croissante :
\[
0 = T_0 < 1 = T_1 < 3 = T_2 < 6 = T_3 < 10 = T_4 < \dots < T_m < T_{m+1} < \dots
\]
Les intervalles d'entiers semi-ouverts $I_k = [\![ T_k, T_{k+1} - 1 ]\!] = \{ m \in \mathbb{N} \mid T_k \le m < T_{k+1} \}$ pour $k \in \mathbb{N}$ forment une partition de $\mathbb{N}$ :
\[
\mathbb{N} = \bigsqcup_{k=0}^{+\infty} [\![ T_k, T_{k+1} - 1 ]\!].
\]
Puisque tout entier $n \ge 1$ appartient à un et un seul ensemble de la partition, il existe un unique $k \in \mathbb{N}$ tel que $T_k \le n < T_{k+1}$.
\end{correctionbox}

% ----------------------------------------------------
% EXERCICE 37
% ----------------------------------------------------
\subsection*{Exercice 37}
On considère l'équation $ax^2 + bx + c = 0$ avec $a \neq 0$ et on suppose que cette équation admet deux solutions $x_1$ et $x_2$.
Exprimer en fonction de $a$, $b$ et $c$ les expressions suivantes :
\begin{enumerate}
    \item $x_1^2 + x_2^2$
    \item $\dfrac{1}{x_1} + \dfrac{1}{x_2}$ \quad (en supposant $c \neq 0$)
    \item $\dfrac{1}{x_1^2} + \dfrac{1}{x_2^2}$ \quad (en supposant $c \neq 0$)
    \item $\dfrac{x_1 + 1}{x_2 + 1} + \dfrac{x_2 + 1}{x_1 + 1}$ \quad (en supposant $x_1 \neq -1$ et $x_2 \neq -1$)
\end{enumerate}

\begin{correctionbox}
\textbf{Rappel des relations fondamentales de Viète :}\\
Puisque $x_1$ et $x_2$ sont les deux racines du trinôme $ax^2 + bx + c$, on a :
\[
S = x_1 + x_2 = -\frac{b}{a} \quad \text{et} \quad P = x_1 x_2 = \frac{c}{a}.
\]

\begin{enumerate}
    \item \textbf{Calcul de $x_1^2 + x_2^2$ :}\\
    À l'aide de l'identité remarquable $(x_1 + x_2)^2 = x_1^2 + x_2^2 + 2x_1 x_2$, on a :
    \[
    x_1^2 + x_2^2 = (x_1 + x_2)^2 - 2x_1 x_2 = S^2 - 2P.
    \]
    En remplaçant $S$ et $P$ par leurs valeurs :
    \[
    x_1^2 + x_2^2 = \left(-\frac{b}{a}\right)^2 - 2\left(\frac{c}{a}\right) = \frac{b^2}{a^2} - \frac{2c}{a} = \boxed{\frac{b^2 - 2ac}{a^2}}
    \]

    \item \textbf{Calcul de $\dfrac{1}{x_1} + \dfrac{1}{x_2}$ :}\\
    Sous l'hypothèse $c \neq 0$ (assurant que $P \neq 0$, donc $x_1 \neq 0$ et $x_2 \neq 0$), on met au même dénominateur :
    \[
    \frac{1}{x_1} + \frac{1}{x_2} = \frac{x_1 + x_2}{x_1 x_2} = \frac{S}{P}.
    \]
    D'où :
    \[
    \frac{1}{x_1} + \frac{1}{x_2} = \frac{-\frac{b}{a}}{\frac{c}{a}} = \boxed{-\frac{b}{c}}
    \]

    \item \textbf{Calcul de $\dfrac{1}{x_1^2} + \dfrac{1}{x_2^2}$ :}\\
    De même, en réduisant au même dénominateur :
    \[
    \frac{1}{x_1^2} + \frac{1}{x_2^2} = \frac{x_1^2 + x_2^2}{(x_1 x_2)^2} = \frac{S^2 - 2P}{P^2}.
    \]
    En injectant le résultat de la question 1 :
    \[
    \frac{1}{x_1^2} + \frac{1}{x_2^2} = \frac{\frac{b^2 - 2ac}{a^2}}{\left(\frac{c}{a}\right)^2} = \frac{\frac{b^2 - 2ac}{a^2}}{\frac{c^2}{a^2}} = \boxed{\frac{b^2 - 2ac}{c^2}}
    \]

    \item \textbf{Calcul de $\dfrac{x_1 + 1}{x_2 + 1} + \dfrac{x_2 + 1}{x_1 + 1}$ :}\\
    La condition $x_1 \neq -1$ et $x_2 \neq -1$ signifie que $-1$ n'est pas racine du trinôme, soit $a(-1)^2 + b(-1) + c \neq 0 \iff a - b + c \neq 0$.
    Mettons les deux fractions rationnelles au dénominateur commun $(x_1 + 1)(x_2 + 1)$ :
    \[
    E_4 = \frac{(x_1 + 1)^2 + (x_2 + 1)^2}{(x_2 + 1)(x_1 + 1)}.
    \]
    $\bullet$ Développons le numérateur $N$ :
    \[
    \begin{aligned}
    N &= (x_1^2 + 2x_1 + 1) + (x_2^2 + 2x_2 + 1) \\
    &= (x_1^2 + x_2^2) + 2(x_1 + x_2) + 2 \\
    &= (S^2 - 2P) + 2S + 2 \\
    &= \frac{b^2 - 2ac}{a^2} + 2\left(-\frac{b}{a}\right) + 2 \\
    &= \frac{b^2 - 2ac - 2ab + 2a^2}{a^2}.
    \end{aligned}
    \]
    $\bullet$ Développons le dénominateur $D$ :
    \[
    D = x_1 x_2 + (x_1 + x_2) + 1 = P + S + 1 = \frac{c}{a} - \frac{b}{a} + 1 = \frac{a - b + c}{a}.
    \]
    $\bullet$ En faisant le quotient $E_4 = \frac{N}{D}$ :
    \[
    E_4 = \frac{\frac{2a^2 + b^2 - 2ab - 2ac}{a^2}}{\frac{a - b + c}{a}} = \boxed{\frac{2a^2 + b^2 - 2ab - 2ac}{a(a - b + c)}}
    \]
\end{enumerate}
\end{correctionbox}

% ----------------------------------------------------
% EXERCICE 38
% ----------------------------------------------------
\subsection*{Exercice 38}
Un objet de masse $m$ est soumis aux forces gravitationnelles de deux corps de masse $M_1$ et $M_2$. On note $d$ la distance entre les deux corps et $x$ la distance entre l'objet et le corps de masse $M_1$.
\begin{enumerate}
    \item Calculer $x_e$, la valeur de $x$ telle que l'objet soit à l'équilibre.
    \item Montrer que $x_e$ se met sous la forme $x_e = \dfrac{d}{1 + A}$ où $A$ est une fonction qui ne dépend que des masses $M_1$ et $M_2$.
\end{enumerate}
\textbf{Données :} Force gravitationnelle exercée par un corps de masse $M$ sur un objet de masse $m$ :
$F = \mathcal{G}\dfrac{mM}{r^2}$, où $r$ est la distance entre l'objet et le corps et $\mathcal{G}$ la constante de gravitation universelle.

\begin{correctionbox}
\textbf{1. Calcul de la position d'équilibre $x_e$ :}\\
Plaçons les corps sur un axe gradué : le corps de masse $M_1$ est à l'origine ($x = 0$) et le corps de masse $M_2$ à l'abscisse $x = d$. L'objet de masse $m$ est situé entre les deux corps, à la distance $x \in ]0, d[$ de $M_1$.
Sa distance au corps $M_2$ est donc $d - x$.
\begin{itemize}[leftmargin=*,label=\textbullet]
    \item La force gravitationnelle exercée par $M_1$ attire l'objet vers $M_1$ avec une intensité :
    \[
    F_1 = \mathcal{G} \frac{m M_1}{x^2}.
    \]
    \item La force gravitationnelle exercée par $M_2$ attire l'objet vers $M_2$ avec une intensité :
    \[
    F_2 = \mathcal{G} \frac{m M_2}{(d - x)^2}.
    \]
\end{itemize}
Ces deux forces ayant des directions colinéaires et des sens opposés, l'objet est à l'équilibre si et seulement si leurs normes sont égales :
\[
F_1 = F_2 \iff \mathcal{G} \frac{m M_1}{x^2} = \mathcal{G} \frac{m M_2}{(d - x)^2}.
\]
Comme la constante gravitationnelle $\mathcal{G} > 0$ et la masse de l'objet $m > 0$ sont non nulles, on simplifie :
\[
\frac{M_1}{x^2} = \frac{M_2}{(d - x)^2} \iff M_1 (d - x)^2 = M_2 x^2.
\]
Puisque $x \in ]0, d[$, les quantités $x$ et $d - x$ sont strictement positives. En extrayant la racine carrée des deux membres :
\[
\sqrt{M_1}(d - x) = \sqrt{M_2} x.
\]
Développons et isolons l'inconnue $x$ :
\[
\sqrt{M_1} d - \sqrt{M_1} x = \sqrt{M_2} x \iff (\sqrt{M_1} + \sqrt{M_2})x = \sqrt{M_1} d.
\]
On en déduit l'unique position d'équilibre située entre les deux corps :
\[
\boxed{x_e = \frac{\sqrt{M_1}}{\sqrt{M_1} + \sqrt{M_2}} d}
\]

\textbf{2. Mise sous la forme réduite demandée :}\\
Divisons le numérateur et le dénominateur de cette fraction par $\sqrt{M_1} > 0$ :
\[
x_e = \frac{\frac{\sqrt{M_1}}{\sqrt{M_1}}}{\frac{\sqrt{M_1} + \sqrt{M_2}}{\sqrt{M_1}}} d = \frac{1}{1 + \frac{\sqrt{M_2}}{\sqrt{M_1}}} d = \frac{d}{1 + \sqrt{\dfrac{M_2}{M_1}}}.
\]
L'expression est bien de la forme :
\[
\boxed{x_e = \frac{d}{1 + A} \quad \text{avec} \quad A = \sqrt{\frac{M_2}{M_1}}}
\]
La grandeur adimensionnelle $A$ ne dépend exclusivement que du rapport des deux masses $\dfrac{M_2}{M_1}$.
\end{correctionbox}

% ----------------------------------------------------
% EXERCICE 39
% ----------------------------------------------------
\subsection*{Exercice 39}
Deux triangles équilatéraux ont le même centre et leurs côtés sont parallèles deux à deux. L'aire de l'un est égale au double de l'aire de l'autre. Le côté du plus petit triangle a pour longueur $1$.
Calculer la distance $d$ entre les côtés parallèles.

\begin{correctionbox}
\textbf{1. Rapport d'homothétie et côté du grand triangle :}\\
Soient $T_1$ le petit triangle équilatéral et $T_2$ le grand triangle équilatéral.
Notons $c_1 = 1$ la longueur du côté de $T_1$ et $c_2$ celle de $T_2$.
Puisque les triangles ont leurs côtés parallèles et même centre $O$, ils sont homothétiques par rapport à leur centre commun $O$.
Le rapport des aires de deux figures semblables est égal au carré du rapport d'homothétie $k > 0$ :
\[
\frac{\mathcal{A}(T_2)}{\mathcal{A}(T_1)} = k^2 = 2 \implies k = \sqrt{2}.
\]
La longueur du côté du grand triangle équilatéral vaut donc :
\[
c_2 = k \times c_1 = \sqrt{2} \times 1 = \sqrt{2}.
\]

\textbf{2. Calcul de la distance du centre aux côtés (apothème) :}\\
Rappelons que dans un triangle équilatéral de côté $c$, la hauteur $h$ mesure :
\[
h = c \frac{\sqrt{3}}{2}.
\]
Le centre $O$ du triangle équilatéral (qui est à la fois centre de gravité, orthocentre et centre des cercles inscrit et circonscrit) est situé au tiers inférieur de chaque médiane/hauteur en partant de la base.
La distance $r$ du centre $O$ à n'importe quel côté du triangle (rayon du cercle inscrit) est donc :
\[
r = \frac{1}{3} h = \frac{1}{3} \times \left( c \frac{\sqrt{3}}{2} \right) = \frac{c\sqrt{3}}{6}.
\]
Appliquons cette formule aux deux triangles :
\begin{itemize}[leftmargin=*,label=\textbullet]
    \item Distance du centre $O$ aux côtés de $T_1$ :
    \[
    r_1 = \frac{c_1\sqrt{3}}{6} = \frac{\sqrt{3}}{6}.
    \]
    \item Distance du centre $O$ aux côtés de $T_2$ :
    \[
    r_2 = \frac{c_2\sqrt{3}}{6} = \frac{\sqrt{2}\sqrt{3}}{6} = \frac{\sqrt{6}}{6}.
    \]
\end{itemize}

\textbf{3. Calcul de la distance $d$ entre deux côtés parallèles :}\\
Puisque les deux triangles ont la même orientation par rapport à leur centre commun (rapport d'homothétie positif $k = \sqrt{2} > 0$), deux côtés parallèles se situent dans la même direction par rapport au centre $O$.
La distance $d$ séparant ces deux droites parallèles est donc la différence de leurs distances au centre :
\[
d = r_2 - r_1 = \frac{\sqrt{6}}{6} - \frac{\sqrt{3}}{6} = \frac{\sqrt{6} - \sqrt{3}}{6} = \frac{\sqrt{3}(\sqrt{2} - 1)}{6}.
\]

\textbf{Conclusion :}\\
La distance exacte entre les côtés parallèles est :
\[
\boxed{d = \frac{\sqrt{6} - \sqrt{3}}{6} = \frac{\sqrt{3}(\sqrt{2} - 1)}{6}} \approx 0{,}1196
\]
\end{correctionbox}

% ----------------------------------------------------
% EXERCICE 40
% ----------------------------------------------------
\subsection*{Exercice 40}
Soit $ABC$ un triangle rectangle en $A$ tel que $AB = 4$ et $AC = 3$. Calculer la distance minimale d'un point de l'hypoténuse à $A$.

\begin{correctionbox}
\textbf{Méthode 1 : Propriété géométrique de la projection orthogonale}\\
Soit $M$ un point quelconque appartenant au segment $[BC]$ (l'hypoténuse).
La distance de $A$ à $M$ est minimale lorsque le vecteur $\vec{AM}$ est orthogonal à la droite $(BC)$, c'est-à-dire lorsque $M$ coïncide avec le projeté orthogonal de $A$ sur la droite $(BC)$, noté $H$.
En effet, pour tout point $M$ de $[BC]$, le triangle $AHM$ est rectangle en $H$, et d'après le théorème de Pythagore :
\[
AM^2 = AH^2 + HM^2 \ge AH^2 \implies AM \ge AH,
\]
avec égalité si et seulement si $HM = 0$, soit $M = H$.
Puisque les angles en $B$ et $C$ du triangle rectangle sont aigus, le pied de la hauteur $H$ se situe strictement à l'intérieur du segment $[BC]$.
La distance minimale cherchée est donc la hauteur $AH$.

$\bullet$ Calculons l'hypoténuse $BC$ via le théorème de Pythagore dans le triangle $ABC$ rectangle en $A$ :
\[
BC = \sqrt{AB^2 + AC^2} = \sqrt{4^2 + 3^2} = \sqrt{16 + 9} = \sqrt{25} = 5.
\]
$\bullet$ Exprimons l'aire du triangle $ABC$ de deux manières différentes :
\[
\mathcal{A}(ABC) = \frac{AB \times AC}{2} = \frac{BC \times AH}{2}.
\]
D'où :
\[
AB \times AC = BC \times AH \iff AH = \frac{AB \times AC}{BC} = \frac{4 \times 3}{5} = \frac{12}{5} = 2{,}4.
\]

\textbf{Méthode 2 : Optimisation d'un trinôme du second degré (analytique)}\\
Munissons le plan d'un repère orthonormé d'origine $A$, où l'axe des abscisses porte $[AB)$ et l'axe des ordonnées porte $[AC)$.
Dans ce repère, les sommets ont pour coordonnées $A(0\,;\,0)$, $B(4\,;\,0)$ et $C(0\,;\,3)$.
La droite $(BC)$ a pour équation :
\[
\frac{x}{4} + \frac{y}{3} = 1 \iff y = -\frac{3}{4}x + 3.
\]
Un point $M(x, y)$ appartient à l'hypoténuse $[BC]$ si et seulement si $x \in [0, 4]$ et $y = 3 - \frac{3}{4}x$.
Le carré de la distance $AM$ s'exprime comme une fonction de l'abscisse $x$ :
\[
f(x) = AM^2 = x^2 + y^2 = x^2 + \left(3 - \frac{3}{4}x\right)^2.
\]
Développons cette expression :
\[
f(x) = x^2 + \left(9 - \frac{9}{2}x + \frac{9}{16}x^2\right) = \frac{25}{16}x^2 - \frac{9}{2}x + 9.
\]
$f$ est une fonction polynôme du second degré de la forme $ax^2 + bx + c$ avec $a = \frac{25}{16} > 0$.
Elle atteint son minimum absolu en son sommet :
\[
x_0 = -\frac{b}{2a} = \frac{9/2}{2 \times (25/16)} = \frac{9/2}{25/8} = \frac{9}{2} \times \frac{8}{25} = \frac{36}{25} = 1{,}44 \in [0, 4].
\]
La valeur minimale du carré de la distance est alors :
\[
\begin{aligned}
f(x_0) &= \frac{25}{16}\left(\frac{36}{25}\right)^2 - \frac{9}{2}\left(\frac{36}{25}\right) + 9 \\
&= \frac{25}{16} \times \frac{1296}{625} - \frac{162}{25} + 9 \\
&= \frac{81}{25} - \frac{162}{25} + \frac{225}{25} = \frac{144}{25}.
\end{aligned}
\]
On en déduit la distance minimale :
\[
AM_{\min} = \sqrt{\frac{144}{25}} = \frac{12}{5} = 2{,}4.
\]

\textbf{Conclusion :}\\
La distance minimale d'un point de l'hypoténuse à $A$ vaut :
\[
\boxed{d_{\min} = \frac{12}{5} = 2{,}4}
\]
\end{correctionbox}

% ----------------------------------------------------
% EXERCICE 41
% ----------------------------------------------------
\subsection*{Exercice 41}
Une colonne de fourmis de longueur $\ell$ se déplace à vitesse constante $v_c$. La dernière fourmi remonte la colonne jusqu'à sa tête puis rejoint sa place à vitesse constante $v_f$. Pendant cet aller et retour, la colonne a avancé de la longueur $\ell$.
\begin{enumerate}
    \item On note $t_1$ le temps pour rejoindre la tête et $t_2$ le temps pour rejoindre la queue. Exprimer $t_1$ et $t_2$ en fonction de $v_c$ et $v_f$.
    \item Déterminer $\dfrac{v_f}{v_c}$ puis déterminer la distance totale parcourue par la fourmi lors de cet aller et retour.
\end{enumerate}

\begin{correctionbox}
\begin{enumerate}
    \item \textbf{Expression de $t_1$ et $t_2$ (cinématique galiléenne) :}\\
    Puisque la fourmi parvient à remonter la colonne, sa vitesse propre est strictement supérieure à celle de la colonne : $v_f > v_c$.
    Plaçons-nous dans le référentiel en translation rectiligne uniforme lié à la colonne de fourmis (vitesse $v_c$ par rapport au sol) :
    \begin{itemize}[leftmargin=*,label=\textbullet]
        \item \textbf{Trajet aller (de la queue à la tête) :}
        La fourmi et la colonne progressent dans le même sens.
        La vitesse relative de la fourmi par rapport à la colonne est :
        \[
        v_{\text{rel, aller}} = v_f - v_c.
        \]
        La distance à couvrir dans ce référentiel est la longueur $\ell$ de la colonne.
        La durée nécessaire pour atteindre la tête est donc :
        \[
        \boxed{t_1 = \frac{\ell}{v_f - v_c}}
        \]
        \textit{Remarque dans le référentiel terrestre :} À l'instant $t_1$, la tête s'est déplacée de $v_c t_1$. La fourmi a donc parcouru la distance $\ell + v_c t_1 = v_f t_1 \iff t_1(v_f - v_c) = \ell$, ce qui donne la même relation.

        \item \textbf{Trajet retour (de la tête à la queue) :}
        La fourmi fait demi-tour et se déplace dans le sens opposé au mouvement de la colonne.
        La vitesse relative de rapprochement de la fourmi et de la queue est :
        \[
        v_{\text{rel, retour}} = v_f + v_c.
        \]
        La distance à couvrir pour regagner sa place à l'arrière est de nouveau la longueur $\ell$.
        La durée nécessaire pour retrouver la queue est donc :
        \[
        \boxed{t_2 = \frac{\ell}{v_f + v_c}}
        \]
    \end{itemize}

    \item \textbf{Détermination du rapport des vitesses et de la distance parcourue :}
    \begin{itemize}[leftmargin=*,label=\textbullet]
        \item \textbf{Équation temporelle :}
        La durée totale de l'aller-retour est $T = t_1 + t_2$.
        Pendant cette même durée $T$, la colonne avance à la vitesse constante $v_c$ d'une distance exactement égale à sa longueur $\ell$ :
        \[
        v_c T = \ell \iff T = \frac{\ell}{v_c}.
        \]
        En égalant les deux expressions de la durée totale $T$ :
        \[
        \frac{\ell}{v_f - v_c} + \frac{\ell}{v_f + v_c} = \frac{\ell}{v_c}.
        \]
        Comme la longueur $\ell > 0$ est non nulle, nous simplifions par $\ell$ :
        \[
        \frac{1}{v_f - v_c} + \frac{1}{v_f + v_c} = \frac{1}{v_c}.
        \]

        \item \textbf{Introduction du rapport sans dimension $x = \dfrac{v_f}{v_c}$ :}
        Divisons le numérateur et le dénominateur de chaque terme du membre de gauche par $v_c > 0$ :
        \[
        \frac{1}{\frac{v_f}{v_c} - 1} + \frac{1}{\frac{v_f}{v_c} + 1} = 1.
        \]
        En posant $x = \dfrac{v_f}{v_c}$ (avec la contrainte physique $x > 1$) :
        \[
        \frac{1}{x - 1} + \frac{1}{x + 1} = 1.
        \]
        Réduisons au même dénominateur :
        \[
        \frac{(x + 1) + (x - 1)}{(x - 1)(x + 1)} = 1 \iff \frac{2x}{x^2 - 1} = 1.
        \]
        Comme $x > 1$, on a $x^2 - 1 \neq 0$. En effectuant le produit en croix :
        \[
        x^2 - 1 = 2x \iff x^2 - 2x - 1 = 0.
        \]

        \item \textbf{Résolution de l'équation du second degré :}
        Le discriminant du trinôme est :
        \[
        \Delta = (-2)^2 - 4(1)(-1) = 4 + 4 = 8 = (2\sqrt{2})^2 > 0.
        \]
        Les racines sont :
        \[
        x_1 = \frac{2 - 2\sqrt{2}}{2} = 1 - \sqrt{2} < 0, \qquad x_2 = \frac{2 + 2\sqrt{2}}{2} = 1 + \sqrt{2}.
        \]
        Comme le rapport de vitesses est supérieur à $1$, on retient la seule racine admissible :
        \[
        \boxed{\frac{v_f}{v_c} = 1 + \sqrt{2}}
        \]

        \item \textbf{Distance totale parcourue par la fourmi dans le référentiel terrestre :}
        La fourmi se déplace à la vitesse constante $v_f$ pendant toute la durée totale $T = \dfrac{\ell}{v_c}$.
        La distance totale parcourue par la fourmi par rapport au sol est donc :
        \[
        d_f = v_f \times T = v_f \times \frac{\ell}{v_c} = \left(\frac{v_f}{v_c}\right)\ell.
        \]
        En remplaçant le rapport $\frac{v_f}{v_c}$ par sa valeur exacte :
        \[
        \boxed{d_f = (1 + \sqrt{2})\ell \approx 2{,}414\,\ell}
        \]
    \end{itemize}
\end{enumerate}
\end{correctionbox}

\end{document}
